\chapter{Financing, ownership and the group boundary}
\label{ch:financing-groups}

A company can preserve the same operating assets while changing who supplies
its funding and who shares its results. The following cases move from debt
cash flows to ownership, then to the reporting boundary of a group. Each
numerical example states its own contract; the coupon-bearing convertible
below and the later zero-coupon compound example are different instruments.

\section{Debt and hybrid financing}

The contract layer for a debt instrument stores issuer and holder perspective,
currency, issue and settlement dates, issue proceeds and legal face amount.
Its payment terms include coupon/index/spread, reset and fixing dates,
day-count convention, payment calendar, amortisation schedule and maturity.
Its rights and protections include seniority, collateral, covenants, options,
conversion terms, guarantees and linked hedges. Missing terms are
explicit unknowns. A single carrying amount cannot stand in for this record.

The accounting layer stores the classification decision and its evidence,
initial measurement, carrying amount, accrued interest, effective-rate method,
unamortised fees or discounts, fair-value disclosures, impairment or credit
adjustments where applicable, and presentation mappings. Contractual interest,
effective-interest expense, realised cash payments and fair-value changes are
separate quantities. The zero-coupon example in Chapter~\ref{ch:execution}
provides a minimal conformance case for this separation.

The lifecycle must support issuance, recognition, accrual, reset, coupon,
principal repayment, modification, derecognition, default, recovery and
conversion. A modification event first determines which accounting operation
the selected standard requires; it does not always mean issuing a new loan.
An embedded option can affect valuation and classification without being
exercised. A derivative linked as a hedge requires both the economic contract
and any applicable accounting designation; economic risk reduction alone does
not establish hedge-accounting eligibility.

Hybrid instruments make interpretation particularly consequential. IAS 32's
overview discusses liability/equity classification and the separation of a
compound instrument \citep{ias32}. The executable classification must consider
the full terms and the selected standard version. US GAAP or ASBE treatment
must be implemented under its own rules rather than inferred from an IFRS
classification. Where the terms are not disclosed sufficiently, alternative
classifications can coexist only if each is actually compatible with the
known terms and applicable rules.

\subsection{Apple: debt balances and cash repayment}

Apple's borrowing shows why the distinction between a contract and its carrying
amount matters for a funding forecast. A fall in reported debt might suggest
how much cash the company repaid. Its annual disclosure lets us test that
interpretation by separating principal from the adjustments included in the
carrying amount. The following amounts are in millions of US dollars at the
displayed precision; changes run from 30 September 2023 to 28 September 2024
\citep[Note 9, printed pp.~43--44]{apple2024}.

\begin{center}
\begin{tabular}{lrrr}
\toprule
Term-debt component & 2023 & 2024 & Change\\\midrule
Principal &106,572&97,341&$-9,231$\\
Net premium, discount and issuance costs &$-356$&$-321$&$+35$\\
Hedge accounting fair-value adjustment &$-1,113$&$-358$&$+755$\\\midrule
Carrying amount &105,103&96,662&$-8,441$\\\bottomrule
\end{tabular}
\end{center}

The cash-flow statement reports no net term-debt issuance proceeds and cash
repayments of 9,958 million for fiscal 2024 \citep[printed p.~33]{apple2024}.
To connect that cash movement to the table, first calculate the difference
between the change in reported principal and net term-debt cash flow. It is
$-9{,}231-(-9{,}958)=727$ million. Adding the reported measurement-component
changes then gives the carrying-amount movement. With $\Delta$ denoting closing
less opening balance, the two displayed-amount identities are
\[
\begin{aligned}
 \Delta\text{principal}
   &=-9{,}958+727=-9{,}231,\\
 \Delta\text{carrying amount}
   &=-9{,}958+727+35+755=-8{,}441.
\end{aligned}
\]
The 727 is an unallocated balance-to-cash difference. It does not identify an
exchange movement or a journal. Likewise, the 35 is the change in a net
unamortised balance, rather than an identified issuance fee or interest charge;
the 755 is a change in the debt's hedge adjustment, rather than a cash receipt.
Together they explain 790 of the 1,517 million difference between the
carrying-amount movement and net cash repayment. A forecast that used the
8,441 carrying decline as repayment would therefore compute the wrong cash
amount for this observed year. The cash statement's proceeds measure is net;
it must not automatically become legal face issued in another period.

The maturity schedule answers a different funding question. Its six future
principal buckets sum to 97,341 million, with 10,930 in fiscal 2025. The current
carrying amount is 10,912. Their 18 million difference does not allocate a
discount or hedge adjustment to the current portion. The measurements differ,
and their time boundaries must also be retained: Note 1 supplies the fiscal-year
default, the issuance table explicitly labels its 2024--2062 range as calendar
years, and MD\&A separately reports amounts payable within twelve months.
MD\&A also reports future term-debt interest payments of 38.5 billion dollars,
including 2.6 billion within twelve months \citep[printed pp.~25, 34 and 43]{apple2024}.
Those cash requirements cannot replace period interest expense or a present
value. The dated Level 2 fair-value disclosures likewise measure something
different: the annual note reports 88.4 billion dollars at September 2024,
while the interim note reports 82.7 billion at December 2024 alongside a
rounded carrying amount of 94.8 billion
\citep[Note 6, printed p.~10]{apple2025q1}.

\subsection{Net borrowing can conceal refinancing}

Commercial paper exposes a second limitation of aggregate financing data.
Apple describes unsecured short-term promissory notes used for general
corporate purposes, including dividends and repurchases. The annual note says
their maturities are generally less than nine months; that description does
not identify each note or trace individual borrowed dollars to a use.
At the annual statement amounts, commercial paper increases from 5,985 to
9,967 million, a rise of 3,982, while net cash proceeds are 3,960. The resulting
22 million difference remains unallocated. During the following quarter the
balance falls from 9,967 to 1,995, while net cash repayments are 7,944; the
corresponding unallocated difference is $-28$ million
\citep[Note 9 and balance and cash-flow statements]{apple2024}
\citep[Note 6 and balance and cash-flow statements]{apple2025q1}.

Even a complete explanation of those differences would leave refinancing
activity unresolved. Annual Note 9 reports net cash for paper with maturities
of at most 90 days, but separate proceeds and repayments for longer maturities.
In fiscal 2022, for example, the longer-maturity proceeds of 5,948 and
repayments of 7,257 give net repayments of 1,309 million. Adding short-maturity
net proceeds of 5,264 gives total net proceeds of 3,955. For fiscal 2024 the
longer-maturity cash rows are zero and the short-maturity net inflow is 3,960.
Neither that net amount nor a closing balance reveals how many issues and
redemptions occurred between reporting dates.

Consider a hypothetical incremental borrowing of 1,000 million US dollars,
issued at par and repaid after 60 days. Stipulate simple interest of 5\% using
actual days over a 360-day year, with no fees, taxes, hedges or early settlement.
The rate is an assumption of this example; Apple's disclosed 5.00\% weighted
rate at September 2024 does not establish the rate on such a transaction.
Issuing and repaying the principal adds $1{,}000-1{,}000=0$ to net principal
cash flow and leaves zero principal outstanding, yet supplies 1,000 million
between those dates. The cost over the 60 days is
\[
 \text{interest}=1{,}000\times0.05\times\frac{60}{360}
               =\frac{25}{3}\approx8.3333
               \quad\text{million US dollars}.
\]
Issue debits cash and credits principal 1,000. Accrual debits interest expense
and credits interest payable $25/3$. Settlement debits principal 1,000 and
interest payable $25/3$, and credits cash $3{,}025/3$. Each entry balances;
after settlement the incremental liability is zero, while cash and equity are
each lower by $25/3$ relative to the otherwise funded baseline.

This example isolates what a net principal-flow observation fails to reveal.
It is not a reconstructed Apple history: interest expense, interest cash and
other disclosures could distinguish or exclude the added borrowing. Within
that scope, it shows why a financing model needs dated issuance, accrual and
redemption events. Matching a net financing row does not establish either
intraperiod liquidity or funding cost.

\subsection{A rate scenario needs an exposure definition}

The contract's coupon is also insufficient to describe exposure after hedging.
Apple's annual debt note states that interest-rate swaps effectively convert
some fixed-rate US-dollar notes to floating rates, while cross-currency swaps
effectively convert some foreign-currency notes to US-dollar notes. The
effective-rate ranges include interest, discount or premium amortisation and
applicable hedge adjustments. These are relationships between separate
contracts and measurements, not a single rate that can be applied to total
principal \citep[Note 9, printed p.~43]{apple2024}.

The filing supplies a useful, more specific constraint. Annual Item 7A reports
that a hypothetical increase of 100 basis points at all tenors, including
associated derivatives, would increase annual term-debt interest expense by
139 million dollars at September 2024 and 194 million at September 2023
\citep[printed p.~27]{apple2024}. Multiplying Note 4's 2024 interest-rate
contract notional of 14,575 million by 1\% instead gives 145.75 million. It does
not reproduce the reported sensitivity, even allowing for nearest rounding of
the displayed amounts. The difference cannot be assigned to a particular
reset date or swap feature from these aggregates. For 2023 the same shortcut
gives $19{,}375\times0.01=193.75$, consistent with rounded 194; that numerical
match does not validate the shortcut or recover the issuer's model.

The scenario's measure matters as much as its shock. In the same table, the
2024 investment-portfolio response is a 2,755 million decline in fair value.
It cannot be added to an annual interest-expense increase as though both were
the same loss measure. At December 2024 the reported interest-rate notional
is still 14,575 million, while the carrying amount of term debt subject to
fair-value hedges is a rounded 13.3 billion, compared with 13.5 billion at
September. These aggregates describe different populations or measurements.
The interim statement that market risk has not changed materially refers the
reader to annual Item 7A; it does not publish an updated 139 million sensitivity
\citep[Note 4 and Part I, Item 3]{apple2025q1}.

Hedge accounting and cash presentation also remain distinct. Annual Note 1
presents derivative assets and liabilities at gross fair values. Note 4
classifies derivative cash flows with the items hedged, which are
\emph{generally} operating activities; that qualifier does not classify every
derivative settlement as operating. The maximum disclosed horizon for hedging
debt-related foreign-currency cash-flow variability is 18 years, and Apple
states that some exposures may be left unhedged and that offsets may be partial
\citep[Notes 1 and 4, printed pp.~34 and 37--38]{apple2024}.

Finally, Item 7A's currency-risk statistic has its own population. Apple reports
a Monte Carlo estimate, at 95\% confidence over one day, of 538 million dollars
at September 2024 and 669 million at September 2023 for adverse fair-value
changes in its foreign-currency derivatives. The assessment excludes underlying
currency exposures, whose changes generally offset derivative changes. This
derivative VaR is neither a hard maximum loss nor a measure of the company's
net currency risk; the issuer's simulation has not been replicated here.
A financing simulation can retain these dated scenarios as disclosed
constraints, but still needs actual holdings, currencies, payment calendars,
reset terms and hedge relationships to calculate future cash and accounting
consequences. The next question is how financing choices change the quantities
and claims of the company's owners.


\section{Equity and ownership}

Equity has both monetary accounts and security quantities. The security
register records share class, issued and outstanding quantities, treasury
holdings, par or stated value where applicable, voting rights, dividend rights,
redemption features, options and conversion relationships. Monetary equity
accounts include contributed capital, reserves, retained earnings and other
components required by the selected basis. These records evolve together
through issuance, repurchase, cancellation, split, dividend and conversion
events.

A stock split changes quantities and per-share measures without necessarily
changing aggregate monetary equity. A cash dividend affects distributions and
cash rather than current operating expense under the corresponding ordinary
equity treatment. Repurchases and employee share awards require their own
rules. Earnings per share uses a dated quantity schedule and rules for potential
ordinary shares; simply dividing profit by the closing share count is
insufficient when shares change during the period.

Funding simulations must connect the quantity register to market assumptions.
An equity issue at a scenario-dependent price changes cash raised, shares and
dilution. The model must state whether it targets a fixed cash amount, a fixed
share count or a capital ratio. These are different management decisions even
when the historical statement would report the same aggregate proceeds.

\subsection{A share history changes the denominator}
\label{sec:eps-lifecycle}

A company can earn the same profit with different numbers of shares
outstanding during the year. Dividing by the closing count would attribute
the entire year's earnings to shares that may have existed for only a few
months. Basic earnings per share instead relates earnings attributable to
ordinary shareholders to the time those shares were outstanding. Diluted
earnings per share asks a further, conditional question: how would that
measure change if eligible potential shares had become ordinary shares?
Neither calculation exercises an option or converts a bond in the company's
actual accounts.

The following synthetic case uses IAS~33 in the retained original EU
Regulation 2023/1803, including its integral Appendix~A. It develops a method
under that historical text; applying it to a real issuer requires the
issuer's reporting basis, operative version and instrument terms. It is
separate from the Apple disclosures that follow
\citep[IAS~33, paragraphs 2--4, 9--12, 30--36]{eu20231803}.
Assume one fully paid ordinary-share class, one consolidated reporting
population, and a 365-day year from 1 January to 31 December 2025.
Monetary amounts are currency units; share counts are numbers of shares,
so EPS is currency units per share.

The year opens with 1,000 ordinary shares and no treasury holdings.
On 1 April, consideration becomes receivable for 200 new shares at 12
each and is paid that day: cash and contributed equity rise by 2,400.
On 1 July, a two-for-one split doubles every outstanding ordinary share.
On 1 October, the company buys 300 shares for 7 each and holds them in
treasury. Cash and equity fall by 2,100; issued shares remain 2,400,
treasury shares become 300, and outstanding shares become 2,100.
The split contributes no additional resources. IAS~33 therefore restates
earlier share counts onto the new share basis, whereas the resource-changing
issue and repurchase enter the average from their respective dates
\citep[IAS~33, paragraphs 19--21, 26--29]{eu20231803}.

\begin{center}
\begin{tabular}{lrrr}
\toprule
Interval, including first date & Days & Actual outstanding & Split-adjusted\\
\midrule
1 January--1 April & 90 & 1,000 & 2,000\\
1 April--1 July & 91 & 1,200 & 2,400\\
1 July--1 October & 92 & 2,400 & 2,400\\
1 October--1 January 2026 & 92 & 2,100 & 2,100\\
\bottomrule
\end{tabular}
\end{center}

Each interval excludes its end date, so no day is counted twice. The basic
denominator is consequently
\[
 D=\frac{90(2000)+91(2400)+92(2400)+92(2100)}{365}
   =\frac{162480}{73}\approx2225.7534.
\]
The closing count of 2,100 answers a different question. Nor should the
split-adjusted count overwrite the historical register: there really were
1,000 outstanding shares in January. If the comparative year had constant
shares of 1,000 and ordinary earnings of 6,000, its previously calculated
EPS of 6 would be presented as 3 after this split. The same retrospective
principle applies to a split occurring after year-end but before
authorization of the statements. A later cash issue does not receive that
treatment \citep[IAS~33, paragraphs 64--65, 70(d)--71]{eu20231803}.

\subsection{Attribute earnings before testing potential shares}

Suppose consolidated profit is 15,000 from continuing operations and
18,000 in total. Non-controlling interests receive 1,000, all from
continuing operations. A cumulative, nonparticipating preference class,
classified here as equity, has a current-year entitlement of 2,000.
Deduct that entitlement even though it has not been declared. A payment
of 500 of prior-year arrears is not another deduction from this year's
earnings. These are supplied attribution facts; the EPS calculation does
not infer them from cash dividends
\citep[IAS~33, paragraphs 12--14 and A1]{eu20231803}.

Earnings attributable to the parent's ordinary shareholders are thus
$15000-1000-2000=12000$ from continuing operations and
$18000-1000-2000=15000$ in total. Discontinued operations contribute
3,000. Basic continuing EPS is $12000/D\approx5.3914$; basic total EPS
is $15000/D\approx6.7393$. Their difference, approximately 1.3479,
is basic discontinued EPS. Profit already includes the interest and tax
effects described below; adding interest again at this stage would be wrong.

There are three potential-share series. First, cash options have been
outstanding all year: 200 shares at an exercise price of 4 before the
split, and 400 at 2 afterward. Use a supplied average market price of 5
on the post-split basis. Earlier prices must be converted to that same
basis. The options require no future employee services and contain no
performance condition, debt tender, redemption requirement or cash-settlement
alternative.

Assumed exercise would provide $400(2)=800$. At the average market price,
800 corresponds to 160 fairly priced shares. Only the remaining
$400-160=240$ shares enter the diluted denominator; there is no earnings
adjustment. Adding all 400 shares would ignore the assumed proceeds.
IAS~33 describes this as the difference between assumed issue and the
shares represented by proceeds at the average price
\citep[IAS~33, paragraphs 45--47 and A4--A5]{eu20231803}.
The calculation records neither cash of 800 nor an actual share issue.

Second, a convertible debt series is issued on 1 July. Its original terms
permit voluntary conversion into 600 post-split ordinary shares; conversion
is neither mandatory nor subject to a cash alternative. For this case,
supplied issue proceeds of 10,000 divide into a liability of 9,200 and
conversion equity of 800. The recognized effective-interest expense for
the remaining year is 800, with a related tax benefit of 200 and no other
earnings consequences. The original terms pay a coupon of 400 on 1 October
and principal of 9,200 plus a final coupon of 400 on 1 January 2026.
The supplied effective rate is $1/23$ per 92-day period: discounting those
payments gives $400/(24/23)+9600/(24/23)^2=9200$.
Interest of 400 first raises the liability to 9,600; the October coupon
returns it to 9,200, and another 400 of interest gives a year-end liability
of 9,600 before January settlement. The intervening payment explains why
the two interest amounts are equal. These supplied amounts include the
relevant cost and discount effects; substituting a coupon alone need not do so.
Instrument classification and component valuation are separate accounting
inputs, illustrated by the compound-bond example below.

Assumed conversion from issue eliminates the net expense of
$800-200=600$ and adds shares only for the 184 days that the instrument
exists. Its contributions are
\[
 \Delta N_b=600,\qquad
 \Delta D_b=600\frac{184}{365}=\frac{22080}{73},
 \qquad
 \frac{\Delta N_b}{\Delta D_b}=\frac{365}{184}\approx1.9837.
\]
IAS~33 requires the after-tax interest effect and any consequential
earnings changes, together with this issue-date weighting
\citep[IAS~33, paragraphs 33--39, 49--50]{eu20231803}.
For example, a profit-sharing expense caused by removing interest would
also change the numerator; it cannot disappear merely because a coupon
and a tax rate are available.

Third, the preference class can convert into 400 post-split ordinary
shares and has existed throughout the year. Its assumed conversion adds
back the current entitlement of 2,000. Earnings per incremental share
are therefore $2000/400=5$. The class has no participation in residual
earnings; participating rights would require a different attribution
calculation.

\subsection{The order of dilution changes the answer}

Continuing-operations earnings attributable to ordinary shareholders,
12,000 here, control admission. Each series is assessed separately in
ascending order of earnings per incremental share: options at zero,
debt at about 1.9837, then preference shares at 5. Recalculate after each
admitted series; comparing every instrument only with basic EPS loses
this dependency \citep[IAS~33, paragraphs 41--44]{eu20231803}.

The options reduce continuing EPS from approximately 5.3914 to 4.8667.
The debt then reduces it further:
\[
 D_{\mathrm{dil}}=D+240+\frac{22080}{73}
                 =\frac{202080}{73},
 \qquad
 E_{\mathrm{dil,cont}}
   =\frac{12000+600}{D_{\mathrm{dil}}}
   =\frac{7665}{1684}\approx4.5517.
\]
The preference class would increase this result because its incremental
ratio, 5, exceeds 4.5517. Exclude it. Yet $5<5.3914$, so it would have
appeared dilutive if tested against basic EPS alone. That is why the
order matters.

The admitted population also determines the other EPS amounts. Total
earnings increase from 15,000 to 15,600 through the same debt addback;
discontinued earnings remain 3,000. Thus
\[
 E_{\mathrm{dil,total}}=\frac{15600}{202080/73}
       =\frac{4745}{842}\approx5.6354,\qquad
 E_{\mathrm{dil,disc}}=\frac{3000}{202080/73}
       =\frac{1825}{1684}\approx1.0837.
\]
Continuing and discontinued EPS sum to total EPS before presentation
rounding. The report retains that reconciliation, each admitted numerator
and denominator increment, and the excluded preference series. Basic and
diluted EPS receive equal prominence; negative EPS is still presented
\citep[IAS~33, paragraphs 66--70]{eu20231803}.

Two sign changes test the control number. If continuing ordinary earnings
are instead a loss of 5,000 while total ordinary earnings remain a profit
of 15,000, none of these three series is admitted: each would reduce the
loss per share from continuing operations. Total profit cannot reverse
that decision. Conversely, retain continuing earnings of 12,000 but
change total ordinary earnings to a loss of 9,000. Options and debt still
enter, so diluted total EPS uses $-9000+600=-8400$ over the same diluted
denominator. It is less negative than basic total EPS. That is the
consequence of the continuing-operations control, illustrated explicitly
in the standard's Appendix~A
\citep[IAS~33, paragraphs 41--43 and A3]{eu20231803}.

\subsection{Actual conversion, expiry and incomplete terms}

Actual events divide the year differently. Suppose the debt pays its coupon
and then converts under its original terms on 1 October; interest ceases
that day. Before conversion, recognized interest is 400 with a tax benefit
of 100. Paying the coupon reduces the liability to 9,200, which transfers
to ordinary equity; the original conversion-equity component remains in equity.
After conversion, 600 actual shares enter basic EPS for 92 days.
The diluted calculation adds only the preceding 92 days of potential
shares and the preconversion net expense of 300. There are therefore
$600(92+92)/365$ additional share-years in the diluted comparison, with
no overlap. Under the otherwise identical supplied earnings assumptions,
basic ordinary earnings rise by the 300 of avoided later expense, and
the remaining addback of 300 restores the same fully diluted numerator
as before. Treating the potential shares as outstanding for all 184 days
\emph{in addition} to the actual shares would double count
\citep[IAS~33, paragraphs 21(c), 36--38]{eu20231803}.

Expiry has no subsequent actual-share leg. If the options instead lapse
on 1 October, their contribution is $240(273/365)$, using a supplied
average price of 5 for this variant's relevant period. They do not vanish
from the earlier part of the year. An exercise price equal to or above
the average price gives zero incremental shares in this simple option
profile. More generally, a conversion that leaves continuing EPS exactly
unchanged is not dilutive: paragraph 41 requires a decrease in profit
per share or an increase in loss per share.

Employee awards require another input even when their other terms fit
the option calculation. Under paragraph 47A, assumed proceeds include
the IFRS~2 value of future goods or services. If that supplied value were
400 across these 400 options, the assumed price would be
$2+400/400=3$, and the increment at market price 5 would be 160 shares.
This conditional arithmetic neither values the services nor calculates
award expense. Performance-based awards follow the contingent-share
provisions instead; share-or-cash alternatives, written puts, rights
issues and participating or partly paid shares also need their own
methods \citep[IAS~33, paragraphs 47A--63, A2 and A13--A16]{eu20231803}.

An EPS assessment therefore references a dated book, reporting population,
share class, attributed earnings, share history, prices and complete
instrument terms. It produces report amounts and their reconciliation,
without changing the book. Missing average prices, future-service values,
conversion consequences or settlement terms leave diluted EPS unresolved
even when basic EPS is available. Missing attribution or ordinary-share
timing can prevent both. Known complex terms select the relevant further
method; they are not grounds for silently dropping the instrument.
These distinctions let the following issuer disclosures supply facts
without being mistaken for a complete method or a complete history.


\subsection{Apple: shares through time}
\label{sec:apple-eps-ownership}

Apple's disclosures show why the ownership record needs dates as well as
quantities. In fiscal 2024, its opening 15,550,061 thousand common shares,
less 499,372 thousand repurchased and plus 66,097 thousand issued net of
employee-tax withholding, give 15,116,786 thousand closing shares. That bridge
accounts for the change in the reported stock. It does not say how long each
share participated in the year's earnings. The reported basic weighted average
is 15,343,783 thousand shares \citep[Notes 3 and 10, printed pp.~36 and 44]{apple2024}.

The unit conversion matters before any timing inference. Net income is
93,736 million dollars, whereas the denominator is in thousands of shares.
Thus $93{,}736\times1{,}000/15{,}343{,}783\approx6.109$ dollars per share,
consistent with reported basic EPS of \$6.11. Using closing shares instead
gives about \$6.201: it computes a different ratio. Adding the reported
64,312 thousand incremental shares from dilutive awards gives the diluted
denominator of 15,408,095 thousand and EPS of about \$6.084, reported as
\$6.08. These are checks using displayed amounts. Conditional nearest-rounding
intervals also overlap the reported two-decimal EPS; the issuer's tie convention
and undisclosed exact amounts remain unverified.

The first quarter of fiscal 2025 makes the distinction again. Income of
36,330 million divided by 15,081,724 thousand basic shares, with the same
factor of 1,000, gives about \$2.409. The diluted denominator adds 69,141
thousand shares and gives about \$2.398, consistent with the reported \$2.41
and \$2.40 \citep[Note 3, printed p.~7]{apple2025q1}. Those denominators describe
three months; adding them to other quarterly averages would not produce an
annual share count. Period lengths and consistent measurement rules are needed
to combine averages.

Even a weighted average leaves dates unresolved. To see exactly what it
constrains, consider a deliberately simplified common-share model for the
364 days of fiscal 2024. Apple reports a 52-week year ending 28 September
2024 \citep[Note 1, printed p.~34]{apple2024}. Let $S_0$ be the opening number
of ordinary shares, and let $\Delta S_j$ be a signed change after $d_j$ whole
days have elapsed since 1 October 2023. Positive changes issue shares;
negative changes repurchase them. An event then affects the remaining
$364-d_j$ days. Assume no splits or other quantity changes, and aggregate all
net issuance into one event. The average is the opening holding plus each
change weighted by its remaining time:
\begin{equation}
 \overline S=S_0+\sum_j\Delta S_j\frac{364-d_j}{364},
 \qquad 0\leq d_j<364.
 \label{eq:apple-weighted-shares}
\end{equation}
All quantities in this equation are ordinary shares and event days are integers.
An immediate change receives full weight; a change after 363 days receives
only one day's weight. Closing shares depend on the sum of changes, while
equation~\eqref{eq:apple-weighted-shares} also constrains their timing. Neither
relation separately identifies every change and date.

The construction must also respect information outside the annual EPS note.
Apple reports fourth-quarter purchases of 35,697, 42,910 and 33,653 thousand
shares in the intervals 30 June--3 August, 4--31 August and 1--28 September,
respectively \citep[Item 5, printed p.~19]{apple2024}. A history putting all
repurchases in the middle of the year fails that disclosure even if its annual
average is correct. The later interim EPS note also reports a first-quarter
fiscal 2024 basic average of 15,509,763 thousand shares. Both constraints are
included in the following two authored quantity histories.

\begin{center}
\small
\begin{tabular}{@{}lrr@{}}
\toprule
Event & History 1 & History 2 \\
\midrule
First-quarter purchase & $-40{,}298\;(45)$ & $-40{,}298\;(45)$ \\
First-quarter purchase & $-40{,}298\;(46)$ & $-40{,}298\;(46)$ \\
Net issuance & $+66{,}097\;(182)$ & $+66{,}097\;(91)$ \\
Further purchase & $-131{,}136\;(179)$ & $-15{,}643\;(159)$ \\
Further purchase & $-175{,}380\;(180)$ & $-290{,}873\;(160)$ \\
June--August interval purchase & $-35{,}697\;(290)$ & $-35{,}697\;(290)$ \\
August interval purchase & $-42{,}910\;(322)$ & $-42{,}910\;(322)$ \\
September interval purchase & $-33{,}653\;(350)$ & $-33{,}653\;(350)$ \\
\bottomrule
\end{tabular}
\end{center}
The table expresses changes in thousands of shares; parentheses give elapsed
days, not negative signs. Each quantity is multiplied by 1,000 when used in
equation~\eqref{eq:apple-weighted-shares}. Treat the displayed source centres
as exact solely for this illustration. Both histories start at 15,550,061,000
shares, repurchase 499,372,000, issue a net 66,097,000 and end at 15,116,786,000.
Their first 91 days average 15,509,763,000 shares: the two equal early purchases
reduce the average by $40{,}298{,}000(46+45)/91=40{,}298{,}000$.
For the full year, the signed time-weighted changes sum to
$-75{,}085{,}192{,}000$ share-days in each history. Dividing by 364 reduces the
opening count by 206,278,000, giving the reported 15,343,783,000 average.
An independent sum of daily balances gives the same result.

After 100 elapsed days, however, the first history has 15,469,465,000 shares
and the second 15,535,562,000. The earlier issuance makes the difference.
The fourth-quarter purchase dates lie inside their reported intervals, but
none of these illustrative dates is an estimate of Apple's actual transactions.
This establishes non-uniqueness for the stated quantity constraints. Prices,
award-level settlement and monetary journals impose further conditions; the
two examples are not a complete reconstruction of the filing or an exhaustive
family, and no probability attaches to either. A query for shares on an
intermediate date must therefore retain the timing alternatives or obtain
additional evidence.

\subsection{Awards, repurchases and cash measure different movements}

The next quarter shows how another disclosure can narrow an unknown. Note 7
describes about 100 million repurchased shares, while Item 2's interval table
gives 41,627, 32,784 and 25,379 thousand, totalling 99,790 thousand. The balance
sheet reports 15,040,731 thousand closing shares against 15,116,786 thousand
at the preceding year end. At the displayed centres, the quantity bridge requires a net
other increase of $15{,}040{,}731-15{,}116{,}786+99{,}790=23{,}735$ thousand
shares \citep[printed pp.~3, 10 and 20]{apple2025q1}. This narrows the total
without identifying the events composing it. Repurchase authority is different
again: both the note and programme footnote explicitly say that the programme
does not oblige Apple to acquire a minimum amount of shares.

Restricted stock units have their own movement record. In thousands of RSUs,
the fiscal 2024 bridge is $180{,}247+80{,}456-87{,}633-9{,}744=163{,}326$:
opening awards plus grants, less vesting and cancellations. The next quarter
gives $163{,}326+61{,}391-35{,}683-2{,}278=186{,}756$
\citep[Note 11, printed pp.~44--45]{apple2024}
\citep[Note 8, printed p.~10]{apple2025q1}. These gross award movements are
different from net common shares issued or incremental shares in the diluted
EPS calculation. Apple reports that most vested RSUs were net share settled
and that approximately 31 million shares were withheld for employee taxes in
2024. Subtracting that rounded amount from 87.633 million vested RSUs does
not explain the 66.097 million net issuance line in full.

The annual report describes several award types under the 2022 employee
plan and the usual terms of its RSUs. In that reported practice, RSUs generally vest
over four years subject to continued employment, settle one-for-one in common
shares and carry dividend-equivalent rights. ``Generally'' cannot determine
the vesting dates of every grant. The prior plan ceased authorising new grants
in the third quarter of 2022; that does not say its outstanding awards vanished.
Likewise, the approximately 24 million RSUs excluded as antidilutive in fiscal
2023 are a period-specific disclosure, not a universal exclusion rule
\citep[Notes 3 and 11, printed pp.~36 and 44--45]{apple2024}. A future dilution
calculation needs the applicable award terms and EPS rules as well as quantities;
the arithmetic above does not validate that methodology.

Monetary statements impose another set of distinctions. The following amounts
are millions of dollars, with reductions and payments shown as positive
magnitudes \citep[printed pp.~32--33 and 45]{apple2024}
\citep[printed pp.~4--5 and 10]{apple2025q1}.
\begin{center}
\small
\begin{tabular}{@{}lrr@{}}
\toprule
Reported measure & Fiscal 2024 & Fiscal 2025 Q1 \\
\midrule
Repurchase reduction in equity & 95,846 & 23,476 \\
Cash-flow repurchases & 94,949 & 23,606 \\
Employee-tax withholding reduction in equity & 5,605 & 2,993 \\
Cash-flow employee-tax payments & 5,441 & 2,921 \\
Share-based compensation in equity & 12,034 & 3,383 \\
Share-based compensation expense & 11,688 & 3,286 \\
\bottomrule
\end{tabular}
\end{center}
The withholding reductions combine the contributed-capital and retained-earnings
rows: $3{,}993+1{,}612=5{,}605$ annually and $1{,}891+1{,}102=2{,}993$ in the
quarter. These differences remain unallocated here; assigning them to excise
tax, transaction fees or unsettled trades would require further evidence.
The narrative's \$95.0 billion and \$23.3 billion purchase amounts also retain
their own precision and transaction scope. The separate employee-tax settlement
narrative cannot silently replace the cash-flow line.

Finally, grant-date value, vesting-date value, compensation expense and future
unrecognised cost answer different questions. Apple's stated policy recognises
expense on a straight-line basis for its estimate of awards that ultimately
vest \citep[Note 1, printed p.~34]{apple2024}. Its unrecognised RSU cost
increases from \$19.4 billion over a weighted-average 2.4 years to \$29.2 billion
over 2.9 years at the quarter end. Dividing either amount by its weighted period
does not disclose a constant annual expense schedule or a cash liability:
different cohorts have different remaining recognition periods, and future
vesting estimates can matter even under the stated straight-line policy.
The ownership model therefore needs linked award cohorts, recognition periods,
share events and cash settlements. Keeping those dates and measurements distinct
allows an actual later grant or settlement disclosure to narrow the histories
without overwriting the reported EPS or inventing a missing event.

\section{A plan's capacity is different from an employee's award}
\label{sec:apple-stock-plan}

The share register has already separated awards, vesting and common shares
issued. The documents governing an award explain why those quantities can move
differently. Apple's 4 March 2022 Form 8-K supplies the 2022 Employee Stock Plan,
a time-based restricted-stock-unit form and a performance-award form as
Exhibits 10.1--10.3. These are the counterparts of three references in the
annual report. The filing reports shareholder approval and effectiveness on
4 March 2022. The plan also records Board adoption on 9 November 2021 and uses
that earlier date as its defined Effective Date when setting its term. The two
dates perform different jobs; replacing them with one undifferentiated date
would lose information needed to interpret the plan
\citep[Items 5.02 and 5.07; Exhibit 10.1, Sections 1 and 15(a)]{apple2022plan}.

Start with the capacity to make awards. Section 5(a) supplies 510 million
shares plus specified amounts available or recycled from the 2014 plan, subject
to a maximum of 1,274,374,682. That maximum is neither the current unused
capacity nor a promise to issue that many common shares. Section 5(b) counts
a full-value award at two reserve shares for each underlying share, whereas
exercised options and stock appreciation rights use one reserve share for each
share to which the exercise relates. Options and stock appreciation rights
count their gross underlying amount even when fewer shares are delivered
\citep[Exhibit 10.1, Section 5(a)--(b)]{apple2022plan}.

A small hypothetical award makes the distinction concrete. Suppose 100
restricted stock units were initially charged against the reserve at two for
one, all 100 vest and 20 shares are withheld for employee taxes. Assume no
other adjustments. The employee receives 80 shares. Under the plan's
full-value withholding rule, the 20 withheld shares return 40 reserve shares,
because the original charge used that ratio. The two quantities are therefore
\[
 \underbrace{100-20=80}_{\text{common shares delivered}},\qquad
 \underbrace{2(100)-2(20)=160}_{\text{net reserve shares used}}.
\]
The reserve records authority used, rather than shares in the employee's
account. An option exercise covering 100 shares can also deliver 80 after
withholding, yet uses 100 reserve shares under the gross-count rule; the
option withholding does not restore capacity. The same observed delivery
therefore cannot identify the reserve movement without the award type and
counting rule. Cash settlement creates another distinction: the plan excludes
the shares that would otherwise have been delivered from its share count.
These are consequences of Section 5(b), not an assumed employee tax rate,
cash payment or compensation expense \citep[Exhibit 10.1, Section 5(b)]{apple2022plan}.

The 2025 proxy illustrates why the reporting population must travel with each
number. At 28 September 2024 its equity-compensation table reports 125,973,077
shares subject to outstanding awards under the 2022 plan, 37,387,521 under the 2014 plan and
10,612 under the director plan, totalling 163,371,210. The 2014 amount includes
39,984 options and 37,347,537 restricted stock units. Removing those options
leaves 163,331,226 shares subject to restricted stock units across the three listed populations.
The annual note's 163,326 thousand restricted stock units is a different
reported amount; these disclosures alone do not identify a reconciliation.
The proxy separately reports 1,022,428,528 shares available under the 2022
plan, 66,924,428 under the employee purchase plan and 4,110,122 under the
director plan, totalling 1,093,463,078. That total includes different plans
from the outstanding-award total and is governed by their counting conventions
\citep[printed p.~97 and table footnotes]{apple2025proxy}. A reconstruction
must retain these populations rather than subtracting unlike totals to infer
grants, dilution or unused authority.

Capacity still does not tell the employee what will be delivered or when.
Section 9 permits restricted-stock-unit agreements to differ: the agreement
specifies vesting, settlement may be in shares, cash or a combination, and
delivery may be deferred under its terms. Dividend-equivalent rights may be
granted; the units themselves carry no voting rights. The filed time-based
form then makes a more specific promise. Its Section 7 provides for shares
equal to vested units less the applicable tax-related withholding, delivered
promptly and no later than two and one-half months after vesting, subject to
the stated conditions. The plan's wider permission does not replace that
form's narrower promise \citep[Exhibit 10.1, Section 9; Exhibit 10.2, Sections
7 and 11]{apple2022plan}.

The form's grant notice leaves the recipient, quantity, award date and vesting
schedule blank. Those blanks matter: the annual report's general four-year
practice cannot fill every individual notice. For example, the 2025 proxy
describes named executives' 1 October 2023 time-based awards as vesting in
three equal annual instalments beginning 1 April 2026, over approximately
four and one-half years. It also describes termination terms under which
earned or vested rights and delivery dates must be kept separate: Cook's
retirement treatment does not accelerate the originally scheduled delivery
of shares or dividend equivalents
\citep[Exhibit 10.2, grant notice and Section 8]{apple2022plan}
\citep[printed pp.~51 and 53--54]{apple2025proxy}. These examples establish
that a general vesting convention cannot substitute for the applicable award;
they do not identify every employee's terms or a complete amendment history.

For the simulation, the consequence is concrete. Retain the plan version and
reserve rule alongside the actual award terms, vesting conditions, units
earned, settlement medium, delivery date and withholding choice. A permitted
plan alternative becomes a branch for a particular award only when its terms
or explicitly stated unresolved facts allow it. Otherwise cash-settlement
paths would be invented for a share-only promise, or vesting would create
shares before their contractual delivery. Tax amounts, accounting expense
and diluted EPS require their own facts and methods. The performance-award
form poses the next question: how are units earned when the contractual
performance schedule itself has not been filled in?

\section{Performance conditions determine units before settlement}
\label{sec:apple-performance-award}

The performance form turns the distinction between a plan and a grant into a
calculation problem. Its notice leaves the target quantity, grant date,
performance period and vesting schedule blank. Exhibit A likewise leaves the
percentile thresholds and payout percentages blank, although it specifies
interpolation between thresholds, a cap when Apple's absolute return is
negative, and rounding to whole units. The form therefore supplies part of
the method without supplying an individual award's inputs
\citep[Exhibit 10.3, notice and Exhibit A]{apple2022plan}.

The 2025 proxy provides a concrete schedule for the named executives'
1 October 2023 performance awards. It describes a three-year performance
period from the beginning of fiscal 2024 through the end of fiscal 2026, with
vesting on 1 October 2026 subject to the agreements. This is a dated
description of that grant population, not the result of the performance
period. The proxy separately reports results for awards granted on
27 September 2020 and vested on 1 October 2023. Those older results must not
be attached to the newer grants merely because the same proxy reports both
\citep[printed pp.~51--53]{apple2025proxy}.

For the October 2023 schedule, first suppose that the relative-return
percentile has been supplied correctly. Write it as $p$ on a scale from 0 to
100, so that $p=70$ means the seventieth percentile. Let $m(p)$ be the number
of units earned per target unit before applying the negative-return cap and
employment conditions. The proxy gives zero below the twenty-fifth percentile,
25\% of target at that percentile, 100\% at the fifty-fifth and 200\% at the
eighty-fifth or above. Between those points it requires straight-line
interpolation \citep[printed p.~52]{apple2025proxy}. In these units,
\[
 m(p)=\begin{cases}
 0, & p<25,\\[4pt]
 \dfrac14+\dfrac{p-25}{40}, & 25\leq p<55,\\[7pt]
 1+\dfrac{p-55}{30}, & 55\leq p<85,\\[7pt]
 2, & 85\leq p\leq100.
 \end{cases}
\]
The first slope is $(1-1/4)/(55-25)=1/40$; the second is
$(2-1)/(85-55)=1/30$. Thus a supplied percentile of 40 earns 0.625 units per
target unit, while 70 earns 1.5. The threshold at 25 is a jump, not a line
running from zero at some unreported lower percentile. If two unresolved
rank calculations produced only $p=24.9$ or $p=25$, a hypothetical target of
1,000 units would yield either zero or 250 units before other conditions.
Replacing those two possibilities by every quantity from zero to 250 would
create unsupported outcomes.

Relative performance is only one input. Let $R$ be Apple's absolute total
shareholder return over the same performance period, expressed as a fraction.
The source caps vesting at 100\% of target whenever $R<0$, regardless of
relative rank. For this part of the calculation the multiplier is
\[
 g(p,R)=\begin{cases}\min\{m(p),1\},&R<0,\\m(p),&R\geq0.\end{cases}
\]
At the supplied percentile 70, a hypothetical target of 1,000 units therefore
gives 1,500 units when the return is nonnegative but only 1,000 when it is
negative, assuming the other vesting conditions are met. A company can rank
well in a falling market: the rank does not identify the sign of its return.
The example avoids fractional units; Exhibit A's nearest-whole-unit rule still
needs an explicit treatment of exact half-unit ties before implementing all
inputs \citep[Exhibit 10.3, Exhibit A]{apple2022plan}.

Computing $R$ introduces its own source requirements. The form uses
twenty-trading-day beginning and ending average prices, dividend reinvestment,
and adjustments for specified capital changes. Consider a deliberately simple
path with no dividends, distributions or corporate actions. Suppose each of
the twenty beginning prices is 100; nineteen ending-window prices are 80 and
the final price is 120. The beginning average is 100 and the ending average is
$(19\times80+120)/20=82$, giving $R=(82-100)/100=-18\%$. Substituting only
the first and last prices would instead give $+20\%$. With the supplied
percentile held at 70, that substitution would wrongly remove the cap and
increase the hypothetical award from 1,000 to 1,500 units. Even in this simple
case, the observation window changes a contractual consequence
\citep[Exhibit 10.3, Exhibit A, Beginning Price and Ending Price]{apple2022plan}.

The general return and rank calculation requires more than those two averages.
The form defines the comparison group using initial S\&P~500 membership and
continued listing, with specific bankruptcy, successor and merger provisions.
It ranks Apple ahead of a comparison member with equal return. That tie rule
does not by itself specify a unique numerical percentile formula; a library's
default rank conversion cannot silently supply one. Dividends enter the return
through reinvestment, noncash distributions need values, and corporate actions
need the stated adjustments. A calculation must document how those provisions
are implemented, including avoiding a second addition of dividends to a price
series that already includes them. Unresolved ambiguities belong with the
Committee's determination under the form; they are not permission to invent
historical inputs \citep[Exhibit 10.3, Exhibit A]{apple2022plan}.

Employment conditions then determine which otherwise earned units survive.
Under Section 8 of the filed 2022 performance form, ordinary termination
generally cancels unvested rights. Death or disability before vesting instead
allows a fraction of the performance-qualified award to vest on the specified
vesting date. The fraction is the total performance-period days less the days
remaining at severance, divided by total days. For an expressly hypothetical
eligible fraction of $3/5$ and the uncapped 1,500-unit result, 900 units survive
to that date; if the negative-return cap first reduces the result to 1,000,
600 survive. These calculations use the filed form's terms and assume no
other modification, leave adjustment or recoupment event
\citep[Exhibit 10.3, Sections 3--4 and 7--9]{apple2022plan}.

The proxy's termination summary must remain alongside that versioned form,
rather than overwrite it. It describes pro-rata \emph{accelerated} vesting for
named executives other than Cook upon death or disability, and different
provisions for Cook. The 2022 form's Section 8 instead points to the vesting
date. An applicable executed agreement, amendment or further disclosure is
needed to establish which timing governs a particular executive. This is an
unresolved applicability and timing distinction between the inspected sources,
not an established contradiction for one identified award
\citep[printed p.~53]{apple2025proxy}.

Dividend-equivalent rights add a cash consequence to the same conditional
award. The filed form credits ordinary cash dividends using the outstanding
target units and pays the vested proportion in cash, subject to withholding
and the other terms. Suppose one dividend of \$0.25 occurs while all 1,000
target units are outstanding before severance, and there are no other
dividends or adjustments. The initial credit is \$250. A 1.5 performance
multiplier with full eligibility pays \$375 gross; under the assumed $3/5$
eligible fraction, the vested proportion is $900/1{,}000$, so \$225 is payable.
Applying the negative-return cap reduces that conditional payment to \$150.
This cash right is different from the hypothetical dividend reinvestment used
to measure shareholder return; adding it to that return would count a
different entitlement \citep[Exhibit 10.3, Section 5(b)]{apple2022plan}.

The model consequently needs distinct stages: determine the applicable award
and observations; calculate the return and rank under evidenced conventions;
apply the schedule and cap; determine surviving rights; and apply settlement
and withholding terms. The checked examples cover selected consequences with
explicit inputs. They do not recover the October 2026 outcome, reproduce all
individual awards, value compensation expense or establish diluted EPS. Those
questions require further data and their own accounting methods.

\section{The applicable agreement changes which rights survive}
\label{sec:apple-award-versions}

The timing question can be narrowed by reading the later forms already
incorporated into Apple's filings. The annual exhibit list identifies a CEO
performance form effective 25 September 2022, filed as Exhibit 10.2 with the
quarter ended 31 December 2022. It also includes general and CEO performance
forms effective 29 September 2024 as Exhibits 10.20 and 10.22. These are
different versions and participant categories, even though their notices use
similar headings and leave the individual grant particulars blank
\citep[annual exhibit list; Exhibits 10.20 and 10.22]{apple2024}
\citep[Exhibit 10.2]{apple2023q1awards}.

The 2024 general form does not remove the earlier timing distinction.
Section 8 still provides that qualifying death or disability before vesting
allows the performance-qualified award, multiplied by the severance fraction,
to vest on the Vesting Date. Section 7 requires share delivery as soon as
administratively practical thereafter, within two and one-half months and
subject to its conditions. Thus neither inspected general form supplies a
rule to deliver the shares immediately on severance. The proxy's description
of accelerated vesting remains a summary whose application to a particular
award needs further evidence \citep[Exhibit 10.20, Sections 7--8]{apple2024}
\citep[printed p.~53]{apple2025proxy}. Later filing alone is not evidence
that an earlier grant was amended.

The CEO forms change the surviving entitlement. In both the September 2022
and September 2024 versions, Section 8(b) allows unvested units to remain
eligible for the Vesting Date following death, disability, or qualifying
retirement, subject to the Committee's stated authority. These clauses do not
apply the general form's severance fraction. Retirement combines an age of
at least 60 and at least ten years of service with the stated service scope;
the retirement exception normally also requires termination on or after the
first anniversary of the Award Date. Earlier retirement does not receive
that exception unless the Committee determines otherwise. Continued
eligibility still depends on performance and does not make the eventual
number of shares known on departure
\citep[Exhibit 10.2, Sections 3--4 and 8]{apple2023q1awards}
\citep[Exhibit 10.22, Sections 3--4 and 8]{apple2024}.

To isolate the effect, reuse the hypothetical 1,000-unit target and a
performance multiplier of 1.5, with no modification, recoupment or other
condition changing the result. Suppose death or disability occurs at a point
giving a general-form severance fraction of $3/5$. Under the general form,
$1{,}000\times1.5\times3/5=900$ units remain eligible to vest. Under the
CEO form's stated continuation provision, all $1{,}000\times1.5=1{,}500$
units remain eligible. With a negative-return cap reducing the multiplier to
one, the corresponding quantities are 600 and 1,000. These comparisons
hold the performance result fixed to expose the contract difference; they
are not facts about any executive's award.

The difference carries into cash rights without creating an earlier payment.
With the same hypothetical \$250 dividend-equivalent credit established
before departure, the uncapped gross payments are \$225 under the
general-form fraction and \$375 under the CEO continuation provision. The
units and cash rights remain linked to the same applicable terms and
performance conditions. The CEO forms' Section 7 anchors delivery to the
Vesting Date, and Section 5(b) ties dividend equivalents to the award's
vesting and payment terms. Moving those amounts to the severance date would
change the specified cash and share path
\citep[Exhibit 10.2, Sections 5(b), 7--8]{apple2023q1awards}
\citep[Exhibits 10.20 and 10.22, Sections 5(b), 7--8]{apple2024}.

A second boundary is independent of the performance calculation. An employee
aged 60 with ten years of service has not necessarily satisfied the CEO
form's first-award-anniversary condition. Conversely, passing that anniversary
does not establish the age and service conditions. The record therefore needs
the applicable agreement, Award Date, termination date and reason, service
history and any relevant Committee decision. Unknown values remain unknown;
they cannot be converted into automatic eligibility or forfeiture.

The later forms retain unfilled performance schedules and the same kinds of
return, comparison-group and adjustment dependencies. They do not fill the
rank convention or establish the actual October 2023 award terms by
themselves \citep[Exhibits 10.20 and 10.22, Exhibit A]{apple2024}. The
appropriate result of this comparison is a narrower conditional model: the
general and CEO forms support different survival rules, both preserve later
delivery, and the proxy's timing description still needs an award-specific
explanation. A source-reference match establishes which form was retained;
applicability establishes which form can govern a financial path.


\section{A compound bond exposes classification dependencies}

Consider an issuer receiving 100 for a two-year bond paying annual coupons of
5 and principal 100 at maturity, with a holder option to convert that principal
into ten ordinary shares. Assume its currency is the issuer's functional
currency, no other options, fees or contingent settlements, and that the
conversion feature qualifies as equity under the selected rule. A comparable
liability without conversion yields 10\%. The inspected IAS~32-equivalent
text requires separate liability and equity components and measures equity
as the residual \citep[paragraphs 28--32, AG30--AG32]{aasb132}.

The liability's present value is
\[
 L_0=\frac{5}{1.1}+\frac{105}{1.1^2}
     =\frac{11050}{121}\approx91.3223,
 \qquad E_0=100-L_0\approx8.6777.
\]
Issue debits cash 100 and credits these two components. Under the stated
effective-interest schedule, the first coupon leaves liability
$1.1L_0-5=1050/11\approx95.4545$; after the second coupon it is 100.
Redemption debits liability and credits cash 100. Alternatively, conversion
at maturity debits liability and credits equity 100 while issuing ten shares;
the original equity component remains in equity. AG32 supports the absence
of a gain or loss on that maturity conversion. Exact rational valuations are
kept before the declared monetary rounding.

A variable-number share settlement, cash alternative or other clause can
change classification. The contract record must therefore retain these terms;
it cannot select liability/equity accounts independently of the coupon,
conversion and settlement rules. Early repurchase and induced conversion
invoke separate operations: the inspected AG33--AG35 text describes allocation
and the induced-conversion consideration difference. They are not implemented
by changing the maturity date of the simple example. US-GAAP and ASBE
classification require their own dated authority; the IFRS-oriented split is
never silently reused under another basis.

\hypertarget{schedules-and-initial-classification}{%
\section{Schedules and initial classification}\label{method-financing-methods-schedules-and-initial-classification}}

The following liability and ownership methods use selected dated AASB 9 and AASB 132 provisions \citep{aasb9dated,aasb132}. Issuer, holder, contractual cash flow and accounting measurement remain separate throughout the schedule.

DebtContract contains parties, currency, principal, issue/settlement dates, payment calendar, fixed/index/spread/reset terms, compounding/day-count, business-day adjustment, principal changes, fees, options, security, priority and legal default/remedy terms. Each term has evidence/admission provenance. The economic schedule emits due cash independently of the recognition method. A market-value adapter may value that schedule but cannot choose amortised cost, FVOCI, FVTPL or liability/equity classification.

An investor's classification evaluates business model and contractual cash-flow features; a debt label alone cannot establish the AASB 9.4.1.2/4.1.2A conditions. At initial recognition the selected 5.1.1 rule measures fair value and qualifying costs as appropriate. An unknown day-one observable-value condition yields unresolved under 5.1.1A/B5.1.2A. The retained acquisition cases distinguish FVTPL expensed fees, amortised-cost fees, original EIR and POCI credit-adjusted EIR. No rate method inferred from market value can overwrite the contract rate.

For an issuer with a plain amortised-cost debt, net proceeds 100 and one legally due payment 121 at the end of two annual periods, the exact EIR solves 100=121/(1+r)\^{}2, hence r=10\%. Issue debit cash 100/credit debt 100; year-one debit interest 10/credit debt 10; year-two debit interest 11/credit debt 11; redemption debit debt 121/credit cash 121. Outstanding contractual payment stays 121 until settled even while carrying value is 100 or 110. Total interest 21 equals cash outflow 121 minus proceeds 100. An immediate transaction fee 2 changes net proceeds to 98 and the EIR root; retaining 10\% is wrong. The numerical root method is in the numerical-method discussion; an irrational rate is not silently an exact rational.

A fractional period uses the declared convention. Annual-effective 21\% over half a year produces factor 1.1 and interest 10 on 100; simple proration gives 10.5 and is wrong for that contract. A floating coupon resets from the admitted fixing available at the reset date, not a later observed rate. Actual day-count and holiday-calendar versions are immutable inputs. Unknown fixings preserve conditional paths, rather than defaulting to the latest value.

\hypertarget{modification-extinguishment-and-default}{%
\section{Modification, extinguishment and default}\label{method-financing-methods-modification-extinguishment-and-default}}

At a modification, retain original and replacement terms, counterparty, fees, effective date, legal-release evidence and original EIR. Determine derecognition before measurement. AASB 9.3.3.1--3.3.3 and B3.3.1--B3.3.6 govern the selected issuer-liability cases. The quantitative ten-percent test compares the PV of new terms, including specified borrower/lender fees, at original EIR against remaining original PV. It is not a universal asset test, nor proof that every change below ten percent is immaterial; changed currency, rights and other qualitative terms need the applicable analysis. Insufficient evidence remains unresolved. Third-party payment without legal release does not extinguish debt.

On the first anniversary above, old remaining 121 has PV 110. Replace with 133.1 due at the original maturity date, one year after the modification, with no fees and an explicit supported substantial-terms conclusion: PV 121 differs by exactly 10\%, satisfying the inclusive test. Assume separately supported new fair value 118; debit old debt 110, debit extinguishment loss 8, credit new debt 118. The new EIR is 133.1/118\ensuremath{-}1, not automatically the old rate. If the replacement is instead a stated non-substantial same-currency reduction to 115.5, no fees or other changes, PV 105 leads to debit debt 5/credit modification gain 5 under the B5.4.6 recalculation mechanism, then interest 10.5 and settlement 115.5. The actual change below ten percent still needs the qualitative finding. Repurchase of carrying amount 110 for 108 extinguishes the repurchased obligation: debit debt 110/credit cash 108/credit gain 2. Unamortised fees belong in the carrying amount before that comparison, not a second gain adjustment.

If cash available is 80 against debt 121 due, record the due obligation and an unpaid/default economic state; payment 80 leaves contractual unpaid 41 only when that partial payment actually occurs. Do not cancel the remaining liability, recognize a waiver or issue a loan unless a separately admitted event authorizes it. Covenant acceleration, cross-default, cure periods, collateral realization and restructuring require the actual legal terms and priority waterfall. A scenario may explore specified alternatives with explicit assumptions; it may not silently select the cheapest settlement as an accounting fact.

\hypertarget{equity-and-compound-instruments}{%
\section{Equity and compound instruments}\label{method-financing-methods-equity-and-compound-instruments}}

EquityInstrument retains class, voting/distribution/liquidation rights, issued, treasury and outstanding counts, par/legal-capital requirements and conversion terms. Share counts are integer quantities separate from currency balances. Issue 100 ordinary shares at 2 with no fees: debit cash 200/credit equity 200, issued and outstanding 100. Repurchase 10 at 3 for treasury: debit treasury-equity 30/ credit cash 30; outstanding 90, issued 100, treasury 10. Reissue 5 at 4: debit cash 20/ credit treasury-equity 15/credit other equity 5; outstanding 95, treasury 5. There is no trading profit under the selected AASB 132.33 rule; a statutory par split requires that jurisdiction's facts. Dividend declaration creates a liability only at the admitted point at which the obligation exists; proposal of a dividend is not itself payment or automatically an expense.

For a qualifying fixed-for-fixed issuer conversion option in functional currency, AASB 132.28--32 separates liability fair value and residual equity. Assume total issue proceeds 100, standalone debt fair value 90, zero coupon, maturity two years and cash redemption 100 or fixed 10 ordinary shares; no extra derivative, tax or fee. Issue debit cash 100/credit debt 90/credit conversion-equity 10. Effective interest accretes debt 90 to 100. At maturity conversion under original terms debits debt 100 and credits issued-share equity 100; the original equity 10 stays in equity (a permitted internal transfer is not income). Shares increase 10 and no redemption cash leaves. Redemption instead debits debt 100/credit cash 100; unexercised equity is not remeasured into profit. AG 32 supports no conversion gain/loss in that stated case. Variable share count, cash alternatives, foreign currency and changed conversion terms require classification/modification analysis; they cannot be forced into this example.

Financial statement mappings separate interest, modification/extinguishment results, owner transactions and cash settlement. A financing-liability note reconciles opening carrying amount, cash proceeds/repayments, effective interest, fees, conversions, modifications and FX. Total shares and earnings-per-share denominators need dates and dilution conditions; outstanding shares at year end are not a substitute for a weighted period denominator. Full option valuation, hedge accounting, EPS and legal-capital regimes remain explicitly additional methods; the proposed records preserve their required inputs.


\section{Changed conversion terms change the accounting question}
\label{sec:financing-changes}

A maturity conversion, an early repurchase and an inducement to convert all
remove or alter debt, but they do not allocate the resulting consideration in
the same way. The issuer must retain the old conversion rights, the revised
rights, the date of amendment and the actual settlement. A record that says
only \emph{converted} loses the facts needed to distinguish these events.
The following cases use the retained Australian standards for explicit
calendar-2024/2025 for-profit Tier~1 illustrations, without optional later
amendments. They are pretax: tax consequences remain uncomputed.

For an early repurchase with unchanged conversion privileges, the issuer
allocates the consideration between the liability and equity components using
the method of the original separation. The liability allocation uses its
current fair value; the difference from its carrying amount enters profit or
loss, while the equity allocation stays in equity
\citep[AG33--34 and IE39--46]{aasb132}. The standard's example has a
face amount of 1,000, half-yearly coupons of 50 and ten payments remaining.
The original half-year discount rate is 5.5\%; the contemporaneous comparable
nonconvertible debt rate is 4\%. For either rate $r$, the remaining debt value is
\begin{equation}
 V(r)=50\frac{1-(1+r)^{-10}}{r}+1{,}000(1+r)^{-10}.
 \label{eq:changed-financing-pv}
\end{equation}
The first term values the ten coupons; the second values principal at the
last payment. The source rounds the carrying and current fair values to
962 and 1,081. Repurchase consideration of 1,700 then allocates 619 to equity
and produces debt settlement expense of 119. The entries debit debt 962,
expense 119 and equity 619, and credit cash 1,700. Expensing the whole
$1{,}700-962=738$ would misclassify the equity purchase. The original equity
component of 60 remains within equity; it does not cap the 619 deduction.
Exact present-value fractions and the source's rounded presentation remain
separate in the accompanying calculations.

An inducement asks a different valuation question. At the amendment date,
compare the fair value of what the holder receives under the revised
conversion terms with what the holder would have received under the original
terms. Their difference is an inducement loss
\citep[AG35 and IE47--50]{aasb132}. In the retained example, changing the
conversion price from 25 to 20 on face amount 1,000 increases shares from
40 to 50. At the amendment-date share quotation of 40, the additional ten
shares are worth 400. A later quotation of 50 would give 500 and would use
the wrong measurement date.

For an explicitly immediate, unconditional share-only conversion, suppose the
liability's admitted carrying amount is 960 and the original equity component
is 60. This carrying amount is a separate illustrative input. The resulting
entry debits debt 960 and inducement expense 400, and credits issued equity
1,360; 50 shares are issued. Original conversion equity remains in equity.
The 400 expense offsets the corresponding part of the issued equity in total
equity, while derecognizing the debt increases net equity by 960. There is no
cash movement. Noncash financing disclosure preserves the event without
inventing a financing receipt \citep[43--44]{aasb107dated}.

This immediate case does not determine the recognition timing of a delayed,
contingent or cash inducement. Those terms must first establish the obligation.
Nor does an ability to settle in shares establish equity classification.
A derivative with a net-cash alternative fails the all-equity-alternatives
condition, and an obligation to deliver a variable number of shares worth a
fixed sum does not represent a residual ownership interest
\citep[21 and 26--27]{aasb132}. The system must refuse the fixed-for-fixed
method when these facts are missing or contradicted, while retaining the
instrument and its unresolved classification.

\section{A hedge reserve is neither cash nor a second gain}
\label{sec:hedge-lifecycle}

Suppose a company hedges a highly probable inventory purchase. A derivative
gain can arrive before the inventory affects earnings. Cash-flow hedge
accounting preserves that timing through a separate equity reserve, but only
after the relationship qualifies. The designation must identify an eligible
instrument and item, the hedged risk, the risk-management objective and the
assessment method at inception. An economic relationship, the absence of
credit-risk dominance and an appropriate actual hedge ratio are separate
requirements \citep[6.2--6.4 and B6.4.1--19]{aasb9dated}. Favorable price
correlation alone does not establish them. Unknown designation or valuation
information remains unresolved; it cannot be repaired by assigning zero or
placing a gain in other comprehensive income.

Take one documented, offsetting exposure and a derivative worth zero at
designation. The illustration includes the whole derivative, no excluded
option or forward component and no intermediate reserve release. Let $D_t$
be its cumulative gain, including settled gains, and $H_t$ the cumulative
change in the hedged item's present value attributable to the designated
risk. Their signs refer to their economic values. With admitted offsetting
changes, the cumulative effective gain $E_t$ and ineffectiveness $I_t$ are
\begin{equation}
 E_t=\operatorname{sgn}(D_t)\min\{|D_t|,|H_t|\},
 \qquad I_t=D_t-E_t.
 \label{eq:hedge-cumulative-lower}
\end{equation}
The reserve follows the lower absolute cumulative amount; effective changes
enter other comprehensive income and the remainder enters profit or loss
\citep[6.5.11]{aasb9dated}. This formula handles the stated offsetting case;
a complex exposure outside it needs its own measurement method. In particular,
one same-direction period is not itself proof that no economic relationship
exists. Item valuation must include time value and only the item's relevant
features; a hypothetical derivative cannot import features that exist solely
in the actual hedging instrument \citep[B6.4.5 and B6.5.4--6]{aasb9dated}.

Initially, supplied values $D_1=12$ and $H_1=-10$ produce effective gain 10
and profit 2. At the next measurement, $D_2=15$ and $H_2=-14$ produce
effective gain 14 and cumulative profit 1. The second period therefore records
a derivative increase of 3, other comprehensive income of 4 and a loss of 1.
Applying the minimum only to that period's changes would incorrectly give
effective gain 3. The retained cumulative values are essential economic memory.

Now buy the inventory for 114 and settle the derivative for cash of 15.
Settlement removes the derivative asset; it does not by itself release the
14 reserve. Because the forecast transaction has produced a nonfinancial
asset, debit the reserve 14 and credit inventory 14. Inventory cost becomes
100. This direct basis adjustment is not an other-comprehensive-income
reclassification \citep[6.5.11(d)(i)]{aasb9dated}. Selling the inventory for
130 subsequently produces revenue 130 and cost of sales 100.

With opening cash and equity of 200 and no other movements, the full sequence
has cumulative pretax profit 31, other comprehensive income 14 and a separate
equity-to-inventory basis transfer of $-14$. Closing cash and pretax equity
are both 231. The equity reconciliation is
\begin{equation}
 200+31+14-14=231,
 \label{eq:hedge-equity-bridge}
\end{equation}
where omitting the direct basis transfer would overstate equity. Gross
operating cash is $-114+15+130=31$: the identified hedge's cash flows follow
the inventory position's classification \citep[16]{aasb107dated}. A matching
terminal cash total does not excuse a wrong interim reserve, profit or basis.

If the hedged payments are interest, the reserve instead enters earnings when
the interest affects earnings \citep[6.5.11(d)(ii)]{aasb9dated}. With an
effective gain of 8, derivative settlement of 8 and cash interest of 18,
release 8 from the reserve into earnings. Expense is then 10. The release
has an OCI reclassification of $-8$, but no second cash receipt. The
illustrative nonfinancial company consistently elects operating classification
for interest under the retained rule; net operating cash is $-10$
\citep[31--33]{aasb107dated}. A later presentation regime needs its own
applicability decision.

Stopping hedge accounting also requires a reason. If the future cash flows
remain expected, the accumulated reserve remains until the flows occur,
subject to the unrecoverable-loss rule. If they are no longer expected, the
reserve is reclassified immediately to profit or loss. No longer highly
probable is not the same as no longer expected
\citep[6.5.11(d)(iii) and 6.5.12]{aasb9dated}. The worked alternatives keep
a gain reserve of 10 in the first case and release it in the second. A loss
reserve of 10 with an unrecoverable portion of 4 instead releases expense 4
and leaves a loss reserve of 6. Each decision changes a named exposure and
reserve; it cannot be inferred from the derivative's cash settlement alone.

\subsection{Fair-value adjustments and benchmark changes}

An admitted fair-value hedge of amortised-cost debt uses a different route.
Both the derivative change and the risk-attributable debt adjustment enter
profit or loss; the latter changes the debt carrying amount
\citep[6.5.8]{aasb9dated}. For debt initially carried at 100, a derivative
gain of 8 and debt-value increase of 10 give a derivative asset of 8, debt
of 110 and net loss of 2. No cash-flow hedge reserve is created.

Once the debt ceases receiving hedge adjustments, the retained adjustment
must be amortized using a recalculated effective interest rate
\citep[6.5.10]{aasb9dated}. Suppose contractual principal of 100 and coupon
of 21 are payable together one year later. From carrying value 110, the
recalculated rate is $121/110-1=10\%$, giving effective interest expense
11 before settlement. Cash interest remains 21 and principal remains 100;
the difference between cash interest and effective expense amortizes the
10 adjustment. Reusing the old 21\% rate would over-accrete the liability.
Missing contractual dates or payments leave that calculation unresolved.

Benchmark reform is not a general exemption from this analysis. The retained
standard's rate-update expedient applies to necessary reform changes whose
new cash-flow basis is economically equivalent to the old basis. Additional
changes then receive the ordinary modification or derecognition analysis
\citep[5.4.5--9]{aasb9dated}. Comparing one 4\% coupon with 3.8\% plus
a 0.2\% spread verifies one number, not equivalence of all future terms.
The amendment record must retain the legal change, rates, spreads, reset
terms, dates and the supporting equivalence analysis.

A qualifying change to formal hedge designation is neither automatic
discontinuation nor a new relationship. Existing cash-flow reserves follow
the alternative benchmark under the specified conditions rather than being
reset to zero. Other qualifying requirements continue to apply; additional
changes and mixed benchmark groups need their own assessment
\citep[6.9.1--10]{aasb9dated}. The rate-specific 24-month exception for a
separately identifiable risk component likewise requires a supported
expectation and ends on the stated adverse reassessment
\citep[6.9.11--13]{aasb9dated}.

The exact conditional journals and adverse alternatives are retained in
\path{data/accounting-master/finite-method-decisions-20261015/financing/};
the field and event specification is
\path{docs/design/proposal/methods/financing-changes-20261015.md}.
These cases specify the financial consequences of admitted values and
eligibility findings. The next section extends the treatment to excluded
components and rebalancing. Actual valuation, delayed or contingent inducements
and group hedges retain their named method and evidence obligations.
Supplying a market value is useful input; it is not a completed valuation method.


\hypertarget{excluded-hedge-components-and-rebalancing}{%
\section{Excluded hedge components and rebalancing}\label{excluded-hedge-components-and-rebalancing}}

The following conditional methods use the retained AASB~9 compilation for calendar-2024/2025 Australian for-profit Tier~1 illustrations, without optional later amendments. Actual authority admission and component valuation remain separate requirements.

\hypertarget{the-option-premium-pays-for-protection}{%
\subsection{The option premium pays for protection}\label{the-option-premium-pays-for-protection}}

A purchased option can protect an inventory purchase even if it expires with no intrinsic value. Its time value is the price of that protection. When only intrinsic value is designated, applying the earlier whole-derivative method would mix the price of protection with the payoff that offsets the purchase. AASB~9.6.5.15 \citep[6.5.15 and B6.5.29--33]{aasb9dated} separates them. The designated intrinsic-value movement follows the applicable hedge method; the related time-value movement accumulates in a different equity component. The nature of the hedged item determines when that cost affects earnings. A transaction such as an inventory purchase uses a cost attached to that transaction. Protection over a period uses systematic, rational amortisation over the period in which the intrinsic-value hedge adjustment can affect earnings. The life of the option alone does not decide which route applies (B6.5.29--30).

The event record must identify the instrument, designated component, underlying item and risk, critical terms, actual and aligned time values at designation and every measurement date, intrinsic-value changes, separate reserve balances, the transaction or time-period classification, amortisation policy, and each basis transfer or reclassification. Existing intrinsic-value effectiveness requirements remain necessary. Unknown component valuations cannot be replaced with a total derivative value and a zero time value.

Suppose a qualifying inventory hedge begins with cash and equity of 200. Buy an option for 6, entirely time value. Cash is 194 and the option asset is 6; there is no initial expense. The first supplied valuation is intrinsic value 12 and aligned time value 4. The option is worth 16. Its increase of 10 comprises an effective intrinsic gain of 12 and a time-value loss of 2, both in OCI but in separate reserves. At purchase and expiry, intrinsic value is 20 and time value zero. The option is worth 20, the intrinsic reserve is 20 and the cost reserve is negative 6. The last period's change of 4 comprises intrinsic gain 8 and time-value loss 4.

Buy the inventory for 120 and settle the option for 20. Settlement removes the option asset and leaves cash 94. The intrinsic reserve reduces inventory cost by 20. Removing the negative cost reserve adds 6 to inventory cost:

\[120 - 20 + 6 = 106.\]

Both direct transfers affect the asset and equity; neither is a second cash flow or an OCI reclassification. Cash 94 plus inventory 106 equals equity 200 after the transfers. Selling for 130 gives profit 24 and cash 224. The cash bridge is \(200-6-120+20+130=224\). The equity bridge is opening 200, profit 24, cumulative OCI 14, and a direct basis transfer of negative 14. Omitting the premium would report cost 100 and overstate profit by 6. For a forecast sale instead, the accumulated cost enters earnings when the sale does. An unrecoverable portion enters earnings immediately under 6.5.15(b)(iii), rather than waiting for a transaction that cannot recover it.

\hypertarget{protection-over-time-and-imperfectly-matching-terms}{%
\subsection{Protection over time and imperfectly matching terms}\label{protection-over-time-and-imperfectly-matching-terms}}

Consider separately an option that protects an existing inventory holding for two equal periods. Admit aligned initial time value 6 and a rational even allocation over those periods. Intrinsic-value and inventory movements are zero in this illustration, isolating the protection cost. After one period, supplied time value 2 gives an OCI loss of 4. Amortisation of 3 is reclassified to expense, leaving the cost reserve at negative 1. After expiry, the final time-value loss of 2 and second amortisation of 3 leave no reserve; cumulative expense is 6. Immediately expensing each observed time-value fall would instead charge 4 and 2. That is a different timing policy. If the intrinsic-value hedging relationship is discontinued after the first period, 6.5.15(c) requires the remaining net reserve loss of 1 to enter earnings immediately. The remaining option asset of 2 is still an asset; discontinuation does not destroy it.

When critical terms differ, B6.5.32 requires an aligned option valuation using terms that perfectly match the hedged item. The hypothetical option is a measurement input for the related cost, not a replacement of the actual contract. Under B6.5.33(a) \citep[B6.5.32--33]{aasb9dated}, if actual time value exceeds aligned time value at inception, base the cost reserve on aligned changes and put the difference between actual and aligned changes in profit or loss. For initial actual/aligned values 6/4 and later values 3/3, actual loss 3 separates into OCI loss 1 and profit loss 2. At expiry, cumulative losses 6/4 leave OCI loss 4 and profit loss 2, before the required transfer or reclassification.

If actual time value is initially lower, B6.5.33(b) limits the reserve by the lower cumulative change of actual and aligned time value; the remainder of actual changes enters earnings. In a same-direction decline example, initial values 2/3 become 1/2.5. Their cumulative losses are 1 and 0.5, so the cost reserve holds loss 0.5 and profit holds loss 0.5. At expiry, cumulative losses are 2 and 3: the reserve loss becomes 2 and cumulative profit from the mismatch returns to zero. The second period therefore records OCI loss 1.5 and a profit gain of 0.5, while actual value falls 1. Applying a fresh lower-of rule only to each period would lose this reversal. These numbers test declining positive time values; they do not define a signed-value selector for every collar or complex combination. Such cases must retain the source-defined alignment and valuation analysis rather than infer it from this small example.

A zero-cost collar can have zero initial time value and nonzero intermediate time value. B6.5.31 still requires the related changes in OCI. A supplied movement from zero to 1.2 and back to zero produces OCI gain 1.2 and loss 1.2; zero initial premium is not permission to omit both. The terminal transaction cost adjustment and initial-cost amortisation are zero in that example.

\hypertarget{forward-elements-and-currency-basis-spreads}{%
\subsection{Forward elements and currency basis spreads}\label{forward-elements-and-currency-basis-spreads}}

AASB~9.6.5.16 \citep[6.5.16 and B6.5.34--39]{aasb9dated} permits the corresponding treatment when the forward element is excluded from a spot-designated forward, or a foreign-currency basis spread is excluded. The accounting election is explicit; designation of spot alone does not automatically make the excluded component an OCI reserve. Apply the transaction/time-period distinction and aligned-value rules in B6.5.34--39. For the forward element, the inception comparison uses absolute actual and aligned amounts. Retain their signed fair values and the changes used for the reserve; do not compare signed inception amounts as though a negative forward element were necessarily smaller in economic magnitude. B6.5.39 carries the same guidance to currency basis spreads. The exact valuation of unmatched forward elements and collars remains a required method for those applications.

\hypertarget{rebalancing-preserves-the-history-of-the-hedge}{%
\subsection{Rebalancing preserves the history of the hedge}\label{rebalancing-preserves-the-history-of-the-hedge}}

A changed hedge ratio can preserve the original risk-management objective. It does not erase the accounting before the change. Under 6.5.5--6.5.6 and B6.5.7--21 \citep[6.5.5--6.5.6 and B6.5.7--21]{aasb9dated}, first measure and recognise ineffectiveness under the existing designation, then adjust designated quantities prospectively. Distinguish ordinary fluctuations around a valid ratio from a changed relationship that requires a new ratio. A changed objective requires discontinuation of that relationship; a counterparty default is not repaired by changing quantities. Update the documented expected sources of ineffectiveness after rebalancing.

Before a change, cumulative derivative gain 12 and offsetting item loss 10 give reserve 10 and profit 2. Immediately before rebalancing, admitted values 18 and negative 14 instead give reserve 14 and cumulative profit 4. Record the period's derivative gain 6 as reserve gain 4 and profit gain 2 before adjusting the quantities. A reset to equal amounts at the new designation date would hide the accumulated ineffectiveness.

Store each retained, added and removed volume with its designation date and valuation reference. An added item layer begins at the rebalancing date; the earlier layer keeps its old reference. The source's example has 100 tonnes designated at forward price 80 and adds 10 tonnes when the price is 90. In a zero-discount illustrative comparison at price 95, their changes are \(100(95-80)+10(95-90)=1550\). Applying the old price to all 110 tonnes gives 1650, inventing a gain of 100 on the added layer. Actual measurement must use the admitted risk, dates, discounting and contractual terms.

Adding an instrument similarly starts its change at the new designation date without resetting the old instrument. Decreasing designated instrument volume leaves any retained undesignated derivative at fair value through profit or loss unless another qualifying designation applies. For example, removing 10\% of a homogeneous instrument from the designation leaves its total carrying value unchanged; of a subsequent total fair-value gain 10, the undesignated 1 enters profit and the designated 9 follows the continuing hedge's effectiveness test. This split does not assert that all 9 is effective. Decreasing item volume invokes partial-discontinuation rules for the removed part. Whether its forecast cash flows remain expected determines the reserve's later release; removal from designation alone does not mean cancellation.

These procedures extend the existing supplied-valuation method. They leave delayed or contingent inducement recognition, award conditions and settlement choices, broader EPS participation, group/net-position hedges and complex component valuation as named financing work, rather than treating a balanced illustration as the completed financing design.


\section{Control and an elimination book}

The consolidation procedure selects the control graph at each date, aligns
supported accounting policies and presentation currencies, combines controlled
entity books, and posts acquisition, elimination and non-controlling-interest
adjustments in a separate book. A control conclusion needs evidence and a
rule version. Ownership percentage alone is insufficient. Unknown purchase
allocation, currency translation or unmatched intragroup differences remain
named unresolved adjustments rather than unexplained balancing entries.

For a simplified acquisition of a different subsidiary from the goodwill case
that follows, assume control is established, all net assets
are already at the relevant acquisition measurements, no tax or goodwill
arises, and non-controlling interest is measured proportionately. A parent
has cash 60, investment 40 and equity 100; an 80\%-owned subsidiary has cash
30, plant 20 and equity 50. Consolidation debits subsidiary equity 50,
credits investment 40 and credits non-controlling interest 10. Group assets
are cash 90 plus plant 20, and group equity is 100 attributable to the parent
plus 10 to other owners.

A subsequent intragroup loan of 10 shifts cash from parent to subsidiary,
creates a receivable and payable, and is eliminated by debiting payable and
crediting receivable 10. Group totals remain 110. An unmatched loan currency
or settlement date would block that exact elimination. If the parent instead
sells its entire interest for 44 immediately after the initial state and
loses control, the group removes subsidiary cash 30, plant 20 and
non-controlling interest 10, recognises cash consideration 44 and gain 4;
remaining cash/equity is 104. More complex acquisitions and retained interests
require additional measurement operations; this example fixes the structure
and arithmetic without certifying a complete consolidation standard.

\hypertarget{perimeter-and-consolidation-state}{%
\section{Perimeter and consolidation state}\label{method-groups-fx-methods-perimeter-and-consolidation-state}}

The acquisition, control and currency methods below use the inspected IFRS 3, IFRS 10 and IAS 21 provisions in the retained 2023 EU text \citep{eu20231803}. The acquisition with goodwill is a new example; it does not change the preceding zero-goodwill case.

ControlAssessment records investor/investee, relevant activities, rights and their substantive/protective character, variable returns, ability to use power, principal/agent analysis, effective interval and evidence. IFRS 10.5--9 requires power, returns and their connection, with reassessment on changed facts. An 80\% share count is evidence to analyze, not an automatic substitute for control. Unknown control yields separate conditional perimeters and unresolved group results; no entity disappears merely to make reconciliation easy.

PerimeterVersion binds control intervals, ownership classes, report cutoffs, functional/presentation currency, policy-alignment adjustments and the exact source entity books. ConsolidationAdjustment is a distinct group book with matched counterparty/event references; legal-entity journals remain intact. The procedure is: select time-compatible entity books; apply supported policy alignment; translate them; add like balances/movements; eliminate investment against acquisition equity; eliminate intragroup balances, income/cost and unrealised gains; allocate profit/OCI to owner/NCI; reconcile all statements and movement notes. IFRS 10.19--23/B86--B96 support these steps. Reporting-date differences need B92--B93's conditions and intervening adjustments; a generic nearest-date join is invalid.

\hypertarget{acquisition-and-ownership-case}{%
\section{Acquisition and ownership case}\label{method-groups-fx-methods-acquisition-and-ownership-case}}

Assume a business combination, no common-control transaction, parent obtaining control and 80\% of ordinary residual interests for cash 96. Acquiree identifiable assets fair value 160 and liabilities 60 give net 100. There is no prior interest, contingent consideration, tax difference, separately unrecognized intangible, transaction cost or impairment; those absences are synthetic facts. Select proportionate-share NCI under IFRS 3.19: NCI 20 and goodwill 96+20\ensuremath{-}100=16 under .32. In parent books debit investment 96/credit cash 96. In the consolidation overlay, debit acquiree acquisition equity 100 and goodwill 16, credit investment 96 and NCI 20. Only post-acquisition income enters group performance. Asset measurement exceptions, deferred tax and a bargain gain reassessment under .36 would change the calculation; goodwill cannot absorb an unexplained missing balance.

If post-acquisition net profit is 20, all accrued in retained earnings with no dividend/OCI, NCI increases 4 to 24 and parent's attributable profit is 16. Parent buys another 10\% for 15 while retaining control: transfer 12 from NCI (10\% of current net 120), credit cash 15 and debit parent equity 3. Remaining NCI 12; goodwill 16 unchanged; no new profit. IFRS 10.23/B96 makes this an equity transaction. This example has equal ordinary rights; preference rights need their actual allocation and cannot use a simple ownership fraction.

In a separate alternative starting before that additional purchase, sell all 80\% for 130 when net assets 120, goodwill 16, NCI 24 and OCI zero. On loss of control, group gain=130+24\ensuremath{-}(120+16)=18. Derecognize individual assets/liabilities and goodwill, remove NCI, recognize consideration and gain. The separate parent book's investment gain is not added again. If 20\% is retained and 60\% sold for 100 with retained fair value 35, group gain=100+35+24\ensuremath{-}136=23; recognize the retained investment 35. Its subsequent associate/financial-asset method needs separately established significant-influence/classification facts. IFRS 10.B98--B99 also requires the appropriate prior-OCI treatment; this case assumes none.

\hypertarget{intercompany-inventory-and-tax}{%
\section{Intercompany inventory and tax}\label{method-groups-fx-methods-intercompany-inventory-and-tax}}

Assume parent's external inventory cost 60 is sold to its controlled subsidiary for 100 on credit, all still held at report date, tax ignored in the base case. Parent records AR 100/revenue 100 and COGS 60/inventory 60; subsidiary inventory 100/ AP 100. Consolidation debits AP 100/credits AR 100 and debits revenue 100/credits COGS 60 and inventory 40. Group inventory 60, revenue and COGS zero, and no group profit arises. When subsidiary later sells externally for 120, group COGS is 60 and profit 60, not 20; release the previously eliminated 40 in the continuing adjustment roll-forward. Reversing a closing overlay and recreating a new one must not recognize the profit twice.

If the purchaser's independently supported tax base is 100, group carrying amount 60 creates deductible difference 40. At assumed applicable 25\% and with all IAS 12 DTA conditions satisfied, recognize group DTA 10/deferred-tax benefit 10; reverse on outside sale. The original tax expense and cash obligations in the seller's tax jurisdiction remain. Different rates, recoverability or tax-law transfer rules change the answer; tax is not simply eliminated with group profit. For an upstream subsidiary sale, attribute elimination and related tax effects to the appropriate subsidiary/NCI performance; downstream parent's profit elimination does not mechanically reduce NCI by the same fraction.

\hypertarget{functional-currency-versus-group-presentation}{%
\section{Functional currency versus group presentation}\label{method-groups-fx-methods-functional-currency-versus-group-presentation}}

EntityCurrencyState names the functional currency and its evidenced rationale, foreign-currency nominal balances, monetary/nonmonetary classification, historical or fair-value measurement dates, rates and sources. IAS 21.21--23 and .28 require transaction-date translation and closing-rate remeasurement of monetary items; historical-cost nonmonetary assets retain transaction-date rates. A functional currency change is prospective under .35--37. Hyperinflation, lack of exchangeability and hedge accounting require separate applicable methods; do not import the historical .26 text into a later amended period without checking it.

Example: buy goods for foreign 100 at functional 1.2 per foreign unit; debit inventory 120/credit payable 120. Closing rate 1.3 gives payable 130 and FX expense 10, while inventory remains 120 before any NRV test. Pay at 1.25: debit payable 130/ credit cash 125/credit FX gain 5. Cumulative expense 5 matches 125 cash less 120 asset cost. Translating the inventory again to 130 would be wrong for this declared basis.

Group presentation is a different operation under IAS 21.39--41. A subsidiary has opening net assets 100, equity 100, no debt or distributions, annual profit 20 entirely earned at rate 1.25, and closing assets 120. Opening historical rate 1.2 and closing 1.3 give translated assets 156, contributed/opening equity 120 and profit 25; translation OCI 11 reconciles the difference. It is explained as 100\ensuremath{\times}(1.3\ensuremath{-}1.2)+20\ensuremath{\times}(1.3\ensuremath{-}1.25)=11, not a free balancing plug. If ownership is 80\%, parent/NCI shares of profit are 20/5 and OCI 8.8/2.2 under the stated equal rights. Use transaction rates; an average is acceptable only as a justified approximation, not when significant fluctuation invalidates it. Goodwill and fair-value adjustments follow IAS 21.47's foreign-operation treatment, rather than freezing them automatically in the parent's currency.

Group reports retain the translation reserve by entity, ownership and event. Loss of control requires the appropriate reserve reclassification under IAS 21.48 and IFRS 10.B99; a no-control-loss change transfers the relevant attribution in equity under applicable conditions. Cash translation effects reconcile cash separately from operating/investing/financing movements. An intercompany monetary balance's exchange result cannot simply be eliminated as if no currency exposure existed; IAS 21.45 and net-investment .32 require the substantive classification.

\section{A foreign subsidiary from acquisition to disposal}
\label{sec:group-fx-disposal}

The preceding examples establish acquisition, translation and disposal
separately. Their interaction matters because a foreign subsidiary carries
more than translated net assets into a sale: goodwill and accumulated exchange
differences also have owners and dates. A sale can leave total equity correct
while putting the wrong amount in profit. We follow one subsidiary through
these steps using the selected historical IFRS 3, IFRS 10, IAS 21 and IAS 7
provisions in the retained EU text \citep{eu20231803}.

A parent starts with USD 300 of cash and equity. USD is its functional and
presentation currency; the subsidiary's functional currency is EUR. At
acquisition the parent pays EUR 96 for 80\% of equal ordinary interests, when
each euro buys USD 1.20. It obtains substantive power over the relevant
activities, exposure to variable returns and the ability to affect them. The
other owners have no substantive rights that defeat that control. Neither
investment-entity treatment nor a common-control transaction applies.

The acquired business already produces outputs and includes a skilled
organised workforce performing a process critical to continuing them, as
required by the assumed IFRS 3.B12C(a) facts. The optional concentration test
is not elected. Identifiable fair values are EUR 20 of cash, inventory of
140 and a payable of 60. There is no other separately identifiable asset,
tax difference, previous holding, contingent consideration or acquisition cost
in this synthetic case. The workforce itself is subsumed into goodwill under
B37; calling a residual goodwill would not justify omitting an identifiable
customer relationship or other asset in a real acquisition.

Identifiable net assets are \(20+140-60=100\) euros. The parent chooses
proportionate-share measurement for the non-controlling interest (NCI) under
IFRS 3.19. NCI is therefore EUR 20 and goodwill under paragraph 32 is
\(96+20-100=16\) euros. This recognised goodwill belongs to the acquired
parent interest; no NCI goodwill has been recognised. The parent's separate
entry debits investment and credits cash by USD 115.2. The consolidation
adjustment debits acquisition equity 120 and goodwill 19.2, and credits the
investment 115.2 and NCI 24. Immediately after acquisition, group cash is
208.8, inventory 168, goodwill 19.2 and the payable 72. Net assets of 324
equal parent equity 300 plus NCI 24.

The subsidiary then sells all the inventory for EUR 160, receives cash, and
pays its EUR 60 payable. Both events occur at USD 1.25 per euro. The specified
event interval has no further expense, impairment or distribution. In the
subsidiary's books, revenue 160 less cost of sales 140 gives profit 20; cash
closes at \(20+160-60=120\) euros. Translated revenue is 200, cost of sales
175 and profit 25. Parent and NCI shares of profit are 20 and 5. The
translation worksheet translates the opening balances to the event rate
before recording these movements; it does not charge cost 175 against
unadjusted acquisition inventory 168 and leave a temporary negative stock.
These are presentation adjustments, not a remeasurement of historical-cost
inventory in the subsidiary's own functional currency.

At reporting, and at an immediate disposal considered below, the rate is
USD 1.30 per euro. Subsidiary cash translates to 156. IAS 21.47 also treats
the acquisition goodwill as an asset of the foreign operation: its EUR 16
becomes USD 20.8. To explain the currency reserve, translate the opening net
assets and the subsequently earned profit from their respective rates:
\[
  \underbrace{100(1.30-1.20)}_{\text{opening net assets}}
  +\underbrace{20(1.30-1.25)}_{\text{later profit}}=11.
\]
Goodwill adds \(16(1.30-1.20)=1.6\). The total translation gain in other
comprehensive income (OCI) is 12.6, but it is not divided mechanically 80/20.
Only the net-asset gain is shared: parent OCI is \(0.8(11)+1.6=10.4\),
and NCI OCI is \(0.2(11)=2.2\). This follows from the chosen proportionate
NCI measurement and the separate attribution required by IAS 21.41 and
IFRS 10.B94. NCI closes at \(24+5+2.2=31.2\), equal to 20\% of the
subsidiary's identifiable net assets of 156. Parent equity is
\(300+20+10.4=330.4\). Group cash 340.8 plus goodwill 20.8 equals their
combined equity of 361.6.

Cash has its own translation bridge. At acquisition, payment of 115.2 less
cash acquired of 24 gives investing outflow 91.2. The subsequent operating
receipt of 200 less payment of 75 gives operating inflow 125. Exchange
changes increase the translated value of the acquired cash by
\(20(1.30-1.20)=2\) and of the later EUR 100 net receipt by
\(100(1.30-1.25)=5\). IAS 7.25--28 reports that total exchange effect of
7 separately from transaction cash flows. Thus
\(300-91.2+125+7=340.8\). The cash exchange effect is 7; the total
translation OCI is 12.6. They measure different populations.

\subsection{Losing control removes the whole subsidiary}
\label{sec:foreign-disposal-reserves}

Suppose an unanticipated transaction immediately after reporting transfers
control to a buyer. No earlier held-for-sale condition is assumed in this
event slice; a real planned disposal needs that assessment. Compare two
alternatives: sell the whole 80\% interest for USD 169, or sell 60\% for
126.75 and retain 20\% with independently admitted fair value 42.25. Both
transactions remove the seller's substantive power. The equal price per
share is an illustrative valuation assumption, not an accounting requirement.
It makes proceeds plus retained value equal to 169 in each alternative.

IFRS 10.B98 removes all subsidiary assets, liabilities and NCI when control
is lost, and measures any retained investment at fair value. If \(C\) is
cash consideration and \(V\) the retained fair value, the gain before the
currency reserve is reclassified is
\[
  C+V+\underbrace{31.2}_{\text{NCI removed}}
      -\underbrace{156}_{\text{subsidiary cash removed}}
      -\underbrace{20.8}_{\text{goodwill removed}}=23.4.
\]
For the full sale, debit cash 169 and NCI 31.2, credit subsidiary cash 156
and goodwill 20.8, and credit gain 23.4. In the retained-interest alternative,
replace the debit to cash 169 with cash 126.75 and retained investment 42.25.
The subsidiary still holds its own bank account; removal of cash 156 is a
change in the consolidated reporting population, not a payment by the parent.

IAS 21.48--48B and IFRS 10.B99 then reclassify the parent's entire accumulated
exchange reserve of 10.4 into profit: debit that reserve and credit disposal
gain. Total disposal gain is \(23.4+10.4=33.8\). NCI's reserve of 2.2
is already within the NCI balance of 31.2 removed above; it is derecognised
without being recycled to profit. Retaining 20\% does not justify recycling
only the fraction corresponding to the shares sold. Paragraph 48A treats a
loss of control as a disposal even when a minority interest remains.

The two alternatives therefore have equal total ending equity but different
cash and investment positions:

\begin{center}
\begin{tabular}{lrr}
\toprule
USD amount & Sell all 80\% & Sell 60\%, retain 20\% \\
\midrule
Cash consideration & 169 & 126.75 \\
Retained investment & 0 & 42.25 \\
Disposal gain including recycling & 33.8 & 33.8 \\
Net disposal cash flow & 13 & $-29.25$ \\
Closing group cash & 353.8 & 311.55 \\
Closing equity & 353.8 & 353.8 \\
\bottomrule
\end{tabular}
\end{center}

IAS 7.39--42 reports acquisition and loss-of-control cash separately, each
net of the whole subsidiary's cash. Net disposal cash is \(169-156=13\)
for the full sale, but \(126.75-156=-29.25\) when 20\% is retained.
A positive parent receipt can therefore accompany a negative consolidated
disposal cash flow. The retained investment's 42.25 is not cash. The two
cash bridges are \(300-91.2+125+7+13=353.8\) and
\(300-91.2+125+7-29.25=311.55\).

For the whole interval, operating profit 25 plus disposal gain 33.8 gives
profit 58.8, attributed 53.8 to the parent and 5 to NCI. Total OCI is
\(12.6-10.4=2.2\), all attributed to NCI; period OCI does not vanish
merely because closing NCI is zero. The equity roll-forward is
\[
  300+\underbrace{24}_{\text{NCI at acquisition}}+58.8+2.2
      -\underbrace{31.2}_{\text{NCI at disposal}}=353.8.
\]
Parent equity also follows directly from \(300+53.8=353.8\). In the full
sale, the parent's separate cost-based investment gain is
\(169-115.2=53.8\), already equal to its cumulative consolidated profit
in this example. Adding that separate gain to the group result would count
the investment return twice.

Intermediate statements are essential checks. Freezing goodwill at USD 19.2
understates both goodwill and parent OCI by 1.6 before disposal. Yet its
larger pre-recycling disposal gain and smaller recycled reserve cancel,
leaving the same final gain of 33.8. Dividing goodwill translation 80/20
would misattribute 0.32 to NCI. Recycling all 12.6 while still removing NCI
31.2 would overstate profit by 2.2 and leave a spurious negative parent
reserve; recycling only the sold fraction in the retained-interest case
would leave 2.6 after control is lost. Final total equity alone misses these
different failures.

The case uses relative event dates and historical technical support. Actual
valuation, tax, impairment, acquisition exceptions and applicable amendments
remain required inputs. A retained 20\% interest also needs a separate
significant-influence and classification assessment before subsequent
measurement; its percentage does not supply that conclusion. The same
separation of legal-entity balances and group adjustments explains why
eliminating an intragroup loan can leave a currency exposure.


\section{Keeping the group calculation attached to its books}
\label{sec:group-record-bindings}

The disposal example gives the parent a cash receipt while removing cash
from the consolidated balance sheet. A software implementation must know
whose cash each entry describes. If it posts the removal of the subsidiary's
156 to the parent's legal book, it records a payment the parent never made.
If it leaves the subsidiary's cash in the group after control ends, it
overstates the group's assets. The consolidation calculation therefore needs
its own book and exact references to the legal books it reads.

Call the legal entities Parent-P and Subsidiary-S, and the reporting entity
Group-P. Parent-P's functional currency is USD; Subsidiary-S's is EUR.
Their identifiers remain distinct even if two imported records have the same
display name. Each legal-book snapshot identifies its entity, book, currency,
relative date and signed balances. At acquisition the parent's investment
115.2 and cash payment stay in its own USD book. The subsidiary's opening
cash 20, inventory 140 and payable 60 stay in its EUR book. Group-P's
consolidation entries refer to both snapshots without altering either.
This implements the separation between combining balances and eliminating
the parent's investment against subsidiary equity
\citep[IFRS~10.B86--88]{eu20231803}.

The selected perimeter has equally precise dependencies. Its control
assessment records the investor, investee, substantive rights, variable
returns and ability to use power, together with the interval for which
those findings hold. Ordinary ownership of 80\% and proportionate NCI
measurement are separate inputs. The currency record distinguishes the
subsidiary's functional currency from the group's presentation currency.
A rate of 1.30 is usable here only when its units are USD per EUR and its
date is the closing date R. An inverse quote or a later rate is a different
input, even when its label looks familiar. The relative A, O, R and D
calendar defines this synthetic history; it supplies no evidence of when an
actual investor could have known its facts.

The group needs more state than the initial corporate example. Goodwill,
the parent's translation reserve and NCI have distinct balances, while
profit, OCI and cash movements retain their period histories. The proposed
\texttt{GroupForeignOperation/1} profile supplies those fields under its own
schema lock. It preserves the common request and result fields rather than
silently adding accounts to the closed corporate profile. Loading this
extension would still require an implementation to declare support for it;
there is no automatic conversion of a corporate state into a group state.

Follow the existing history through that representation. At A the group
records the acquisition. At O it translates the opening balances before the
subsidiary's operating entries. At R it finishes cash and goodwill
translation and allocates operating profit. A and O are intermediate
worksheet states; they are not presented as fully closed interim reports.
The R state retains cash 340.8, goodwill 20.8, parent translation reserve
10.4 and NCI 31.2, together with profit 25, OCI 12.6 and the separate cash
exchange effect 7. Closing the income accounts does not erase those movements.

Full disposal and disposal with a retained 20\% interest then start from
separate copies of that same immutable R state. Their branch identities
prevent the second alternative from being posted after the first. Each
disposal proposal identifies its opening state, control change, source books,
cash consideration, retained fair value, rates and accounting method.
Its permission to post is conditional on this synthetic profile and exact
proposal. It is not permission to post to an actual company's books.

In the retained-interest alternative, the parent's receipt is 126.75.
The group removes subsidiary cash 156, producing investing cash
$126.75-156=-29.25$. The retained investment 42.25 is a separate noncash
asset. Both disposal alternatives remove the entire NCI balance 31.2,
including its translation reserve 2.2. Only the parent's reserve 10.4 is
recycled into the disposal gain. These are checks on the meaning of the
journal, beyond the fact that its debits equal its credits
\citep[IFRS~10.B98--99; IAS~21.48A--B; IAS~7.39--42]{eu20231803}.

The report reads one branch and one dependency view. In both alternatives,
period profit is 58.8, of which 53.8 belongs to the parent and 5 to NCI.
Period OCI is 2.2, attributable to NCI. Ending NCI is nevertheless zero,
because the disposal removes its accumulated balance. The equity movement
must therefore retain both its introduction and removal:
\[
 300+24+58.8+2.2-31.2=353.8.
\]
Reading only the closing NCI account would lose the subsidiary's contribution
while it was controlled. Reading only final equity would also miss the
difference between full-sale cash 353.8 and retained-interest cash 311.55
plus investment 42.25.

Publication follows the common accounting transaction rule. The request
identifies the group owner and expected opening revision; its business-event
identity includes the branch. Before publishing, an implementation must
validate the source and currency bindings, posting authorization, ordered
journals and independently specified state changes. It then publishes the
closing state, journal, economic delta, receipt and new group head together.
The parent and subsidiary legal-book owners remain untouched.

An acknowledgement can be lost after that publication. Retrying the exact
original request then finds the existing receipt and has no second effect,
even though its original opening revision is now old. If publication never
occurred, the same request can commit once. Changing the request identifier,
budget or expected head changes the fingerprint and cannot masquerade as
that retry. A newly proposed event with an outdated opening is rejected.
The documentary packet preserves both complete acknowledgement histories
and an invalid partial-publication example; later runtime tests must exercise
those failures in the chosen storage system.

A corrected exchange-rate source creates a new evidence view. It makes the
old group's calculation unsuitable for a request using the corrected view,
but does not overwrite the old books or remove the old explanation.
Unknown control or an unknown rate instead returns the missing-input
obligation, with no fabricated zero amount. The companion group-binding
specification supplies exact records and exchanges for these cases.
Actual control, valuation, tax, policy alignment and subsequent accounting
for the retained investment retain their separate evidence requirements.


\clearpage
\section{Eliminating the loan, retaining the currency exposure}
\label{sec:intragroup-fx-example}

An intragroup loan disappears from the consolidated balance sheet, yet its
exchange gain can remain in group profit. The reason becomes visible when
we follow both the loan and the money it financed. IAS~21.45 requires the
exchange result to survive elimination; paragraph~32 determines when it
belongs in other comprehensive income instead
\citep[IAS~21.32--33, .45]{eu20231803}.

Consider a parent whose functional and presentation currency is USD and a
wholly owned subsidiary whose functional currency is EUR. Control and both
functional currencies are established assumptions. At an opening rate of
USD1.20 per EUR, the parent's balances are cash 56, an investment in the
subsidiary carried at historical cost 24, and a EUR100 loan receivable worth
120. Its equity is therefore 200. The subsidiary holds EUR120 in cash,
owes EUR100 to the parent and has EUR20 of paid-in equity. At the opening
rate, its equity matches the parent's investment. Eliminating that
investment and the loan leaves group cash and equity of 200.

We isolate an exchange-rate remeasurement of these established balances,
holding every other measurement input fixed. This is neither the initial
valuation of a new loan nor a complete reporting period. Interest, tax,
impairment, hedges and new transactions are outside this calculation;
IAS~21.24 leaves the underlying carrying amount to the relevant measurement
standards. Thus EUR100 is a given monetary carrying amount, not an assumption
that an interest-free long-term loan must be recognised at its face value
\citep[IAS~21.23--24, .50]{eu20231803}. All following statement amounts are
in USD unless explicitly marked EUR.

Now let the rate rise to USD1.30 per EUR. The parent remeasures its receivable
from 120 to 130: debit the receivable 10 and credit exchange gain 10. The
subsidiary records no exchange gain on its EUR loan in its own EUR books.
For group presentation, its cash becomes $120\times1.30=156$ and its payable
becomes $100\times1.30=130$. Its historical capital remains 24. The remaining
2 is translation OCI, explained by its EUR20 net assets and the rate change:
$20\times(1.30-1.20)=2$
\citep[IAS~21.28, .33, .39--41]{eu20231803}.

Consolidation now debits the translated loan payable 130 and credits the
parent's receivable 130. It also debits the subsidiary's historical capital
24 and credits the parent's investment 24. These entries eliminate the
internal balances; neither entry removes the parent's exchange gain. Group
cash is $56+156=212$, with no external liability in this example. Opening
equity 200, exchange profit 10 and translation OCI 2 reconcile to the same
212. Deleting the profit merely because the receivable was internal would
leave equity short by 10. The group still owns the euros whose dollar value
has risen \citep[IFRS~10.B86; IAS~21.45]{eu20231803}.

\clearpage
The loan's settlement facts can change the presentation of that 10.
In the ordinary-loan case, settlement is planned in the foreseeable future,
so the gain stays in consolidated profit. In a separate conditional case,
assume settlement is neither planned nor likely in the foreseeable future
and the loan forms part of the net investment under IAS~21.15. This is a
substantive condition, not an election created by calling a loan long term;
trade receivables and payables are explicitly excluded. The parent's separate
gain is still 10. In the consolidated statements, debit exchange gain 10
and credit OCI 10, giving total OCI of 12 and no exchange profit. The two
cases in Table~\ref{tab:intragroup-fx} have identical group cash and total
equity but different profit and OCI
\citep[IAS~21.15, .32--33]{eu20231803}.

\begin{table}[!ht]
\centering
\caption{Settlement facts change where the exchange result appears.}
\label{tab:intragroup-fx}
\begin{tabular}{@{}lrr@{}}
\toprule
USD amount & Ordinary loan & Qualifying net-investment loan \\
\midrule
Parent's separate exchange gain & 10 & 10 \\
Consolidated exchange profit & 10 & 0 \\
Consolidated OCI & 2 & 12 \\
Consolidated closing cash and equity & 212 & 212 \\
\bottomrule
\end{tabular}
\end{table}

Neither case includes a cash receipt. The same EUR120 remains in the
subsidiary, so its translated value rises by
$120\times(1.30-1.20)=12$. IAS~7.28 presents this exchange effect separately
from operating, investing and financing cash flows. Group cash reconciles
as opening 200, transaction cash flows zero and exchange effect 12, giving
closing 212 \citep[IAS~7.25--28]{eu20231803}. Adding a further gain of 12
would count the same currency movement twice.

The representation must therefore retain both entities' functional currencies,
the monetary balance, the settlement assessment and the origin of each OCI
component. If the settlement facts are unknown, these conditional presentations
remain separate even when they agree on total equity. The group adjustment
does not overwrite the parent's separate journal, and the OCI history remains
available for the applicable later disposal treatment. An equal consolidated
total cannot establish that the histories answer the same profit, cash-flow
or disposal question.


